\documentclass[12pt, a4paper]{report}

% ── Taal & encoding ──────────────────────────────────────────────────────────
\usepackage[dutch]{babel}
\usepackage[utf8]{inputenc}
\usepackage[T1]{fontenc}

% ── Layout ───────────────────────────────────────────────────────────────────
\usepackage[margin=2.5cm]{geometry}
\usepackage{fancyhdr}
\usepackage{parskip}

% Laat TeX moeilijke regels (lange \texttt-paden) als laatste redmiddel rekken
% i.p.v. over de marge te lopen.
\setlength{\emergencystretch}{3em}

% ── Tabellen ─────────────────────────────────────────────────────────────────
\usepackage{booktabs}
\usepackage{longtable}
\usepackage{array}
\usepackage{multirow}
\usepackage{tabularx}
\usepackage[htt]{hyphenat}  % sta woordafbreking toe in \texttt{} (lange klasse-namen in traceability matrix)

% ── Kleur & code ─────────────────────────────────────────────────────────────
\usepackage{xcolor}
\usepackage{listings}

% ── Afbeeldingen ─────────────────────────────────────────────────────────────
\usepackage{graphicx}

% ── Hyperlinks (laad als laatste) ────────────────────────────────────────────
\usepackage{hyperref}
\hypersetup{
    colorlinks = true,
    linkcolor  = black,
    urlcolor   = blue,
    citecolor  = blue,
    pdftitle   = {Security Audit Rapport --- OpenMRS REST Webservices Module},
    pdfauthor  = {LU2 Groep 17},
}

% ── Avans Hogeschool huisstijlkleur ──────────────────────────────────────────
\definecolor{avansRed}{HTML}{A8001A}
\definecolor{avansDark}{HTML}{1A1A1A}
\definecolor{avansGray}{HTML}{6B6B6B}

% ── RAG-kleuren ──────────────────────────────────────────────────────────────
\definecolor{ragAmber}{HTML}{E67E22}
\definecolor{ragGreen}{HTML}{27AE60}
\definecolor{ragYellow}{HTML}{C9A227}
\newcommand{\rood}[1]{\textcolor{avansRed}{\textbf{#1}}}
\newcommand{\amber}[1]{\textcolor{ragAmber}{\textbf{#1}}}
\newcommand{\geel}[1]{\textcolor{ragYellow}{\textbf{#1}}}
\newcommand{\groen}[1]{\textcolor{ragGreen}{\textbf{#1}}}

% \auditfig{pad}{breedtefactor}: toont de afbeelding, of een nette box als het pad ontbreekt.
\newcommand{\auditfig}[2]{%
    \IfFileExists{#1}%
        {\includegraphics[width=#2\linewidth]{#1}}%
        {\fbox{\parbox[c][3.5cm][c]{#2\linewidth}{\centering\ttfamily\small
            [\,afbeelding ontbreekt nog\,]\\[0.4em]#1}}}%
}

% ── Hoofdstuktitel: "1 — Titel" met rode HR ──────────────────────────────────
\usepackage{titlesec}
\titleformat{\chapter}[block]
    {\normalfont\Huge\bfseries\color{avansDark}}
    {\thechapter\ \textcolor{avansRed}{---}\ }
    {0pt}
    {}
    [\vspace{0.4em}{\color{avansRed}\rule{\linewidth}{2.5pt}}\vspace{0.2em}]
\titlespacing*{\chapter}{0pt}{20pt}{18pt}

% ── Header / footer ──────────────────────────────────────────────────────────
\pagestyle{fancy}
\fancyhf{}
\lhead{\small\color{avansGray} LU2 Groep 17 --- Avans Hogeschool}
\rhead{\small\color{avansGray} Security Audit Rapport}
\rfoot{\small\color{avansGray} Pagina \thepage}
\renewcommand{\headrulewidth}{0.4pt}

% ─────────────────────────────────────────────────────────────────────────────
\begin{document}

% ════════════════════════════════════════════════════════════════════════════
%  TITELPAGINA
% ════════════════════════════════════════════════════════════════════════════
\begin{titlepage}
    \newgeometry{margin=0pt}

    % ── Rode balk bovenaan ───────────────────────────────────────────────────
    \noindent
    \colorbox{avansRed}{%
        \parbox[t][3.5cm]{\paperwidth}{%
            \vfill
            \hspace{2.5cm}{\color{white}\fontsize{11}{14}\selectfont\bfseries
                Avans Hogeschool \quad|\quad
                ATIx IN-B2.4 Softwarearchitectuur \& -kwaliteit \quad|\quad
                LU2 Groep 17}
            \vspace{0.6cm}
        }%
    }

    % ── Hoofdinhoud ─────────────────────────────────────────────────────────
    \vspace{2.5cm}
    \hspace{2.5cm}%
    \begin{minipage}[t]{14cm}

        {\fontsize{36}{42}\selectfont\bfseries\color{avansRed}
            Security Audit Rapport}\\[0.5cm]

        {\fontsize{18}{22}\selectfont\color{avansDark}
            OpenMRS REST Webservices Module}\\[0.2cm]

        {\large\ttfamily\color{avansGray}
            openmrs-module-webservices.rest v3.2.0}

        \vspace{0.7cm}
        {\color{avansRed}\rule{14cm}{1.5pt}}
        \vspace{0.6cm}

        \begin{tabular}{@{}ll}
            \textbf{Normenkader}      & NEN-7510:2024 Deel 2 \\[4pt]
            \textbf{Wetgeving}        & AVG, WGBO, NIS2/Cyberbeveiligingswet \\[4pt]
            \textbf{Academiejaar}     & 2025--2026 P4 \\[4pt]
        \end{tabular}

        \vspace{0.6cm}
        {\color{avansRed}\rule{14cm}{1.5pt}}
        \vspace{0.6cm}

        \begin{tabular}{@{}ll}
            \textbf{Datum}   & \today \\[4pt]
            \textbf{Versie}  & 1.0 \\[4pt]
            \textbf{Status}  & Definitief \\[4pt]
        \end{tabular}

    \end{minipage}

    % ── Rode balk onderaan ───────────────────────────────────────────────────
    \vfill
    \noindent
    \colorbox{avansRed}{%
        \parbox[t][1.2cm]{\paperwidth}{%
            \vfill
            \hspace{2.5cm}{\color{white}\small\bfseries
                Vertrouwelijk --- uitsluitend bestemd voor de opdrachtgever}
            \vspace{0.3cm}
        }%
    }

    \restoregeometry
\end{titlepage}

% ════════════════════════════════════════════════════════════════════════════
%  INHOUDSOPGAVE
% ════════════════════════════════════════════════════════════════════════════
\tableofcontents
\newpage

% ════════════════════════════════════════════════════════════════════════════
%  1. EXECUTIVE SUMMARY
% ════════════════════════════════════════════════════════════════════════════
\chapter{Executive Summary}
\label{ch:executive-summary}

\section{Overall Status}
\label{sec:overall-status}

De OpenMRS REST Webservices Module (v3.2.0) is beoordeeld tegen
NEN-7510:2024-2, het normenkader voor informatiebeveiliging in de zorg. De
module ontsluit medische gegevens en patiëntdossiers via een
programmeerkoppeling; de eisen aan de beveiliging zijn daarom hoog. De audit
combineerde een normgerichte controle, een risico- en dreigingsanalyse,
geautomatiseerde scans van zowel de eigen code als de meegeleverde externe
componenten, en een afsluitende, gerichte aanval (penetratietest) op een
afgeschermde testomgeving.

Dit leverde zes uitgewerkte bevindingen op (B-001 t/m B-006). \textbf{Alle zes
zijn tijdens de audit hersteld en opnieuw getest.} Het ging onder meer om het
onbeperkt kunnen raden van wachtwoorden, het kunnen overnemen van de sessie
van een ingelogde medewerker, het ongewenst kunnen uitlezen van interne
systeem- en accountgegevens, en het kunnen vervalsen van logregels die
forensisch onderzoek moeten ondersteunen. Voor elke herstelmaatregel is bewijs
vastgelegd: een aanpassing in de broncode, een geautomatiseerde test en een
herhaalde aanval die nu wordt geblokkeerd. Ook de eerder gevonden mogelijkheid
om zónder in te loggen de volledige systeemconfiguratie op te vragen --- het
zwaarst beoordeelde risico --- is inmiddels afgeschermd.

\textbf{Overall status: \amber{AMBER (gedeeltelijk)}.} De code van de module
zelf is op de geauditeerde punten op orde, maar volledige naleving is nog niet
bereikt: in de meegeleverde externe softwarecomponenten staan nog bekende
kwetsbaarheden open die om updates vragen, en het afdwingen van versleuteld
dataverkeer ligt op het niveau van de serveromgeving en valt buiten de module.
De directe code-risico's zijn daarmee verholpen; de keten- en
omgevingsmaatregelen vragen nog aandacht.

\section{Top Risico's}
\label{sec:top-risicos}

Onderstaande tabel toont de zes uitgewerkte bevindingen in managementtaal, met
de oorspronkelijke ernst en de actuele status. De ernst is de \emph{inherente}
score vóór mitigatie; alle zes zijn inmiddels hersteld en hertest. De nummering
volgt Hoofdstuk~\ref{ch:bevindingen}.

{\small
\begin{tabularx}{\textwidth}{@{}l X l l@{}}
    \toprule
    \textbf{ID} & \textbf{Omschrijving (managementtaal)} & \textbf{Ernst} & \textbf{Status} \\
    \midrule
    B-001 & Inlogpogingen werden niet beperkt: wachtwoorden raden was onbeperkt
            mogelijk, waardoor een geldig account binnen seconden te kraken was. &
            \rood{20 --- Hoog} & \groen{Gemitigeerd \& hertest} \\
    \addlinespace
    B-002 & Een aanvaller kon na de login van een medewerker diens
            gebruikerssessie overnemen (sessie-fixatie). &
            \rood{15 --- Hoog} & \groen{Gemitigeerd \& hertest} \\
    \addlinespace
    B-003 & Elke ingelogde medewerker, ook met minimale rechten, kon de volledige
            lijst van alle gebruikersaccounts uitlezen. &
            \amber{12 --- Midden} & \groen{Gemitigeerd \& hertest} \\
    \addlinespace
    B-004 & Een gemanipuleerde link kon kwaadaardige code laten uitvoeren in de
            browser van een ingelogde beheerder (XSS). & \amber{12 --- Midden} &
            \groen{Gemitigeerd \& getest} \\
    \addlinespace
    B-005 & Interne foutmeldingen toonden technische systeemdetails aan de
            buitenwereld, wat vervolgaanvallen vergemakkelijkt. &
            \amber{12 --- Midden} & \groen{Gemitigeerd \& hertest} \\
    \addlinespace
    B-006 & Logregels konden door een aanvaller worden vervalst
            (log-injectie), wat forensisch onderzoek kan misleiden. &
            \amber{12 --- Midden} & \groen{Gemitigeerd \& geverifieerd} \\
    \bottomrule
\end{tabularx}
}

\vspace{0.4em}
\noindent Daarnaast is de los daarvan gevonden mogelijkheid om \emph{zonder
inloggen} de volledige systeemconfiguratie uit te lezen (het zwaarst beoordeelde
risico, score 20) afgeschermd. Geen van de uitgewerkte bevindingen staat nog
open. De resterende aandachtspunten liggen niet in de modulecode maar in de
toeleveringsketen en de serveromgeving (zie de roadmap hieronder).

\section{Geprioriteerde Roadmap}
\label{sec:roadmap}

De uitgewerkte bevindingen B-001 t/m B-006 zijn afgerond en daarom niet meer in
de roadmap opgenomen. De onderstaande acties betreffen de resterende
keten- en omgevingsmaatregelen.

{\small
\begin{tabularx}{\textwidth}{@{}l X l@{}}
    \toprule
    \textbf{Prioriteit} & \textbf{Actie} & \textbf{Termijn} \\
    \midrule
    \textbf{Nu} & Bekende kwetsbaarheden in meegeleverde componenten met een
                  veilige update verhelpen: de twee niet-brekende
                  ``quick-win''-updates (\texttt{snakeyaml} naar 2.x,
                  \texttt{spring-web} naar 5.3.34) doorvoeren. & Direct \\
    \addlinespace
    \textbf{Deze sprint} & Versleuteld dataverkeer (TLS) op het niveau van de
                  serveromgeving afdwingen en documenteren (reverse
                  proxy/servlet-container); dekt control A.8.24. & 1 sprint \\
    \addlinespace
    \textbf{Later} & Grote, ingrijpende upgrades van verouderde componenten
                  (Spring~6, Tomcat~9+, vervanging van de verouderde
                  JSON-bibliotheek) gefaseerd inplannen. & Roadmap \\
    \addlinespace
    \textbf{Later} & Het toegangsfilter vervangen door een standaard,
                  declaratief autorisatiekader (Spring Security). & Roadmap \\
    \addlinespace
    \textbf{Later} & De productieschakeling naar het database-account met
                  minimale rechten in een draaiende omgeving verifiëren. & Roadmap \\
    \bottomrule
\end{tabularx}
}

\newpage

% ════════════════════════════════════════════════════════════════════════════
%  2. SCOPE & CONTEXT
% ════════════════════════════════════════════════════════════════════════════
\chapter{Scope en Context}
\label{ch:scope}

\section{Module en Versie}
\label{sec:module-versie}

De geauditeerde software is de \textbf{OpenMRS REST Webservices Module}
(\texttt{org.openmrs.module:webservices.rest}), versie \textbf{3.2.0}.
De module vormt de REST API-laag van het OpenMRS-platform: zij ontsluit
meer dan 70 unieke REST-resources (patiënten, encounters, observaties,
medicatie, e.a.) met bijbehorende sub-resources via HTTP/JSON onder het pad \texttt{/ws/rest/v1/}.

\begin{tabular}{@{}ll}
    \textbf{Artifact-ID}    & \texttt{webservices.rest} \\[4pt]
    \textbf{Versie}         & 3.2.0 \\[4pt]
    \textbf{Taal / runtime} & Java 8+, Spring MVC, Apache Tomcat \\[4pt]
    \textbf{Upstream repo}  & \url{https://github.com/openmrs/openmrs-module-webservices.rest} \\[4pt]
    \textbf{Audit fork}     & \url{https://github.com/GrannyGuard/webservices-rest-audit} \\[4pt]
    \textbf{Audit branch}   & \texttt{main} \\
\end{tabular}

\section{Beoordelingsperiode}
\label{sec:periode}

De audit is uitgevoerd in de periode \textbf{29 mei 2026 t/m 18 juni 2026},
verdeeld over vier sprints conform het projectplan:

\begin{tabular}{@{}lll}
    \textbf{Sprint} & \textbf{Periode} & \textbf{Focus} \\[4pt]
    1 & 29 mei -- 5 juni  & Gap-analyse \& projectorganisatie \\[2pt]
    2 & 5 -- 12 juni      & Risico's, threat modelling \& CI/CD-scanning \\[2pt]
    3 & 12 -- 16 juni     & Attack surface, logging \& code coverage \\[2pt]
    4 & 16 -- 18 juni     & Afronding \& rapportage \\
\end{tabular}

\section{Testomgeving}
\label{sec:testomgeving}

De audit heeft plaatsgevonden op de \textbf{Development- en Test-laag} van de
OTAP-omgeving. Er is geen productiesysteem met echte patiëntdata betrokken.

\begin{tabular}{@{}ll}
    \textbf{OTAP-laag}        & Development / Test \\[4pt]
    \textbf{Codebase}         & Fork \texttt{GrannyGuard/webservices-rest-audit},
                                branch \texttt{main} \\[4pt]
    \textbf{Build}            & Maven 3.x, Java 8, lokale \texttt{mvn verify} \\[4pt]
    \textbf{CI/CD-platform}   & GitHub Actions (gehoste runners, Ubuntu-latest) \\[4pt]
    \textbf{SAST}             & CodeQL \texttt{security-extended}, SonarCloud \\[4pt]
    \textbf{SCA}              & Dependabot, OSV-Scanner, Snyk \\[4pt]
    \textbf{Testdata}         & In-memory H2-database (\texttt{development}-scope);
                                geen herleidbare patiëntdata \\[4pt]
    \textbf{Productie}        & Buiten scope --- geen toegang tot live omgeving \\
\end{tabular}

De scheiding tussen test- en productiedata is geborgd via GitHub Environments
(\texttt{dev}, \texttt{test} en \texttt{prod}): \texttt{prod} kent een approval
gate (vier-ogen, A.8.3) die de release-pipeline daadwerkelijk afdwingt
(\texttt{environment: prod} in \texttt{release-sign.yml}), naast een
protected-branch-policy. Secret- en protection-scheiding voor \texttt{test}/\texttt{dev}
is nog een openstaand verbeterpunt; zie
\texttt{docs/security/02-secure-pipelines/02.md} \S2.

\section{Buiten Scope}
\label{sec:buiten-scope}

De volgende onderdelen zijn expliciet buiten scope gehouden:

\begin{itemize}
    \item \textbf{OpenMRS Core en andere modules}: alleen de REST
          Webservices Module zelf is geauditeerd; de onderliggende kern
          (authenticatie via \texttt{Context}, Hibernate-laag, MariaDB) is
          als vertrouwde component beschouwd tenzij de module er direct mee
          interacteert.
    \item \textbf{Productieomgeving en live patiëntdata}: er is geen
          toegang tot een draaiend productiesysteem; alle tests zijn
          uitgevoerd op de development/test-branch zonder herleidbare
          patiëntgegevens.
    \item \textbf{Infrastructuur en netwerkconfiguratie}: TLS-configuratie,
          reverse proxy, firewalls en serverbeheer vallen buiten de
          module-grens en zijn de verantwoordelijkheid van de deployende
          organisatie.
    \item \textbf{Fysieke beveiliging en organisatorische maatregelen}:
          personeelsbeleid, toegangspassen en fysieke serverruimtes zijn
          niet beoordeeld.
    \item \textbf{Penetratietests op een live productiesysteem met echte
          patiëntdata}: er is geen toegang verkregen tot een productieomgeving.
          De penetratietest (ZAP-scan + handmatige tests, bevindingen F1--F9, zie
          \texttt{05-penetration-tests/PentestReport-OpenMRS.md} en
          Bijlage~\ref{app:pentest}) is wél uitgevoerd op een geïsoleerde
          VirtualBox-testomgeving (Kali Linux 2024, OpenMRS~2.8.7,
          Apache Tomcat~9.0.118) en heeft 9~bevindingen opgeleverd.
\end{itemize}

\section{Normenkader}
\label{sec:normenkader}

Het normenkader voor deze audit is NEN-7510:2024 Deel 2. De norm bevat
aanzienlijk meer controls dan hier behandeld; de onderstaande zeven zijn
geselecteerd op basis van de risico-analyse en de aard van de module: een REST API die bijzondere medische persoonsgegevens verwerkt. Dit zijn
de gebieden waar de risico's het hoogst zijn en waar verbetering de
grootste impact heeft.

De audit richt zich op de volgende zeven controls:

\begin{tabular}{@{}ll}
    \toprule
    \textbf{Control} & \textbf{Onderwerp} \\
    \midrule
    A.8.3  & Toegangsbeveiliging \\
    A.8.5  & Authenticatie \\
    A.8.8  & Beheer van technische kwetsbaarheden \\
    A.8.9  & Configuratiebeheer \\
    A.8.15 & Logging \\
    A.8.25 & Veilige ontwikkelcyclus \\
    A.8.24 & Cryptografie \\
    \bottomrule
\end{tabular}

\vspace{0.8em}
De volledige koppeling tussen deze controls en de traceerbare
bewijsartefacten is opgenomen in de traceability matrix
(zie Bijlage~\ref{app:traceability}).

\section{Toepasselijke Wetgeving}
\label{sec:wetgeving}

\begin{tabularx}{\textwidth}{@{}lX}
    \toprule
    \textbf{Wetgeving} & \textbf{Relevantie} \\
    \midrule
    AVG (GDPR) Art.~9           & Verwerking bijzondere persoonsgegevens (medische data) \\
    WGBO Art.~7:454             & Bewaarplicht en vertrouwelijkheid medisch dossier \\
    NIS2 / Cyberbeveiligingswet & Zorgaanbieders als essentiële entiteit; meldplicht incidenten \\
    CRA (Cyber Resilience Act)  & Beveiligingseisen open-source componenten in de supply chain \\
    \bottomrule
\end{tabularx}

\newpage

% ════════════════════════════════════════════════════════════════════════════
%  3. AUDIT METHODOLOGIE
% ════════════════════════════════════════════════════════════════════════════
\chapter{Audit Methodologie}
\label{ch:methodologie}

\section{Aanpak}
\label{sec:aanpak}

De audit is uitgevoerd in vijf stappen, conform NEN-7510:2024-2 en de
risicobeoordelingsaanpak van ISO/IEC 27005:

\begin{enumerate}
    \item \textbf{Gap-analyse}: per NEN-7510-control is de huidige staat van de
          module beoordeeld op basis van broncode-analyse en configuratieonderzoek.
          Elke gap is geclassificeerd als \emph{aanwezig / gedeeltelijk / afwezig}
          en onderbouwd met concrete codereferenties.
    \item \textbf{Risicoanalyse}: op basis van de gevonden gaps en een CIA/BIV-classificatie
          van de kroonjuwelen zijn risico's gescoord via \emph{Kans × Impact} (schaal 1--5).
          Voor de twee best exploiteerbare hoog-risico-aanvalsketens is een
          bow-tie-analyse uitgewerkt.
    \item \textbf{Threat modelling}: de module is beschreven in C4-diagrammen (context,
          container en component). Op basis daarvan is een STRIDE-analyse uitgevoerd en
          zijn trust boundaries vastgesteld. Gevonden threats zijn gemapt op de
          OWASP API Security Top 10:2023.
    \item \textbf{Verificatie}: bevindingen zijn geverifieerd via SAST (CodeQL),
          SCA (Dependabot, OSV-Scanner, Snyk), logging-tests en een attack surface mapping.
          Gemitigeerde kwetsbaarheden zijn herbeoordeeld met dezelfde tooling.
    \item \textbf{Penetratietest}: een doelgerichte test conform de OWASP Web Security
          Testing Guide~v4.2 is uitgevoerd op een geïsoleerde VirtualBox-testomgeving
          (Kali Linux~2024, OpenMRS~2.8.7, Apache Tomcat~9.0.118). De test omvatte
          authenticatie, sessiebeheer, autorisatie, injectie en serverconfiguratie en
          leverde 9~bevindingen op (zie Bijlage~\ref{app:pentest}).
\end{enumerate}

De werkzaamheden zijn verdeeld over vier sprints (29 mei t/m 18 juni 2026).
Er is uitsluitend gewerkt op de development/test-omgeving; er is geen toegang
geweest tot een productiesysteem met echte patiëntdata.

\section{Tooling}
\label{sec:tooling}

{\small
\begin{tabularx}{\textwidth}{@{}p{2.2cm}p{4.6cm}X@{}}
    \toprule
    \textbf{Categorie} & \textbf{Tool} & \textbf{Doel en inzet} \\
    \midrule
    SAST & CodeQL (\texttt{security-extended}) &
        Statische analyse op elke push/PR; resultaten in GitHub Security-tab
        (SARIF). Bijlage~\ref{app:codeql}. \\
    SAST & SonarCloud &
        Code-kwaliteit en beveiligingsregels; coverage-integratie via JaCoCo.
        Bijlage~\ref{app:screenshots}. \\
    SCA & GitHub Dependabot &
        Primaire kwetsbaarhedenstream op basis van de GitHub Advisory Database;
        264 alerts, herleid tot 75 unieke CVE's. Bijlage~\ref{app:dependabot}. \\
    SCA & OSV-Scanner &
        Onafhankelijke kruisvalidatie op de CycloneDX SBOM via OSV.dev;
        bevestigt 64 van de 75 Dependabot-CVE's. \\
    SCA & Snyk &
        Derde scanner (eigen vulnerability-database); 36 aanvullende CVE's
        boven op Dependabot. Bijlage~\ref{app:snyk}. \\
    SBOM & CycloneDX Maven Plugin &
        Genereert \texttt{sbom.cdx.json} (238 componenten) bij elke push naar
        \texttt{main}. Bijlage~\ref{app:sbom}. \\
    Coverage & JaCoCo &
        Instruction-coverage op de beveiligingsrelevante klassen
        (\texttt{AuthorizationFilter}, \texttt{SessionController1\_9});
        rapport als CI-artefact. \\
    Threat modelling & diagrams.net (\texttt{.drawio}) &
        C4-diagrammen op context-, container- en componentniveau. \\
    Methodiek & STRIDE &
        Systematische dreigingsanalyse per C4-component; 18 threats
        geïdentificeerd. \\
    Normatief kader & NEN-7510:2024-2, OWASP API Top 10:2023 &
        Referentiekader voor controls, dreigingsclassificatie en
        auditcriteria. \\
    \midrule
    Pentest & Burp Suite~CE, Hydra~v9.6, Nmap~7.95 &
        Respectievelijk: request proxy \& inspectie, credential brute
        force, port/service scan. \\
    Pentest & Nikto, OWASP~ZAP, SQLMap &
        Geautomatiseerde webscanner, DAST spider + active scan,
        SQL-injectiescan. \\
    Pentest & ffuf~v2.1.0, custom Python~3.x &
        Endpoint discovery / fuzzing; rate-limiting-test en
        gebruikersenumeratie-script. \\
    \bottomrule
\end{tabularx}
}% end \small

\newpage

% ════════════════════════════════════════════════════════════════════════════
%  4. RISICO-ANALYSE & BEVINDINGEN
% ════════════════════════════════════════════════════════════════════════════
\chapter{Risico-analyse en Bevindingen}
\label{ch:bevindingen}

De risicoanalyse brengt de gevonden gaps, de CIA/BIV-classificatie van de
kroonjuwelen en de STRIDE-dreigingsanalyse samen tot één gescoorde risicomatrix
(\S\ref{sec:risicomatrix}). Daaruit volgen de zes uitgewerkte bevindingen
(\S\ref{sec:bevindingen}). De volledige analyse (de zeven gap-tabellen per
control, de C4-diagrammen, de STRIDE-tabel met achttien threats en de twintig
daaruit gescoorde risico's) is vastgelegd in
\texttt{docs/security/01-security-audit/01.md}; dit rapport vult die matrix aan
met drie empirisch bevestigde pentest-risico's tot drieëntwintig
(\S\ref{sec:risicomatrix}).

\section{Risicomatrix}
\label{sec:risicomatrix}

Risico's zijn gescoord volgens \textbf{Risicoscore = Kans $\times$ Impact}, elk op
een schaal van 1--5 (ISO/IEC 27005). De \emph{risicobereidheid} is laag omdat de
module bijzondere persoonsgegevens (AVG Art.~9) verwerkt: elk risico met score
$\geq 10$ dat de vertrouwelijkheid of integriteit van een patiënt-, medisch- of
authenticatie-kroonjuweel (KJ1, KJ2, KJ3) raakt, geldt automatisch als
\emph{onacceptabel} en vereist directe actie.

\begin{minipage}[t]{0.56\textwidth}
\small
\textbf{Score-grid (Impact $\times$ Kans):}\\[0.4em]
\begin{tabular}{@{}c|ccccc@{}}
 & \textbf{K1} & \textbf{K2} & \textbf{K3} & \textbf{K4} & \textbf{K5} \\
\midrule
\textbf{I5} & \geel{5}  & \amber{10} & \rood{15} & \rood{20} & \rood{25} \\
\textbf{I4} & \groen{4} & \geel{8}   & \amber{12}& \rood{16} & \rood{20} \\
\textbf{I3} & \groen{3} & \geel{6}   & \amber{9} & \amber{12}& \rood{15} \\
\textbf{I2} & \groen{2} & \groen{4}  & \geel{6}  & \geel{8}  & \amber{10}\\
\textbf{I1} & \groen{1} & \groen{2}  & \groen{3} & \groen{4} & \geel{5}  \\
\end{tabular}
\end{minipage}%
\hfill
\begin{minipage}[t]{0.40\textwidth}
\small
\textbf{Classificatie:}\\[0.4em]
\begin{tabular}{@{}ll@{}}
\groen{Laag}        & 1--4 \\
\geel{Laag-Midden}  & 5--8 \\
\amber{Midden}      & 9--12 \\
\rood{Hoog}         & 13--25 \\
\end{tabular}\\[0.6em]
Grenswaarden: $\geq 9$ niet acceptabel
(mitigatie $\leq$ 1 maand); $\geq 13$ onacceptabel
(directe escalatie).
\end{minipage}

\vspace{1em}
De analyse leverde \textbf{drieëntwintig gescoorde risico's} op uit vijf bronnen:
SonarQube (SQ), CI/CD (CD), CodeQL (CQL), het threat model (AT) en de
penetratietest (PT). Eenentwintig daarvan vallen in Midden of hoger; veertien zijn
\emph{onacceptabel} (zeven maal \rood{Hoog} plus zeven Midden-risico's die
KJ1/KJ2/KJ3 met score $\geq 10$ raken). De volgende tabel toont de top van de
matrix (score $\geq 12$); de volledige drieëntwintig staan in
Bijlage~\ref{app:risicomatrix} (de twintig niet-pentest-risico's ook in
\texttt{01.md} \S4.2; de pentestbevindingen in Bijlage~\ref{app:pentest}).

{\small
\begin{longtable}{@{}
    >{\raggedright\arraybackslash}p{1.3cm}
    >{\raggedright\arraybackslash}p{6.2cm}
    >{\centering\arraybackslash}p{1.0cm}
    >{\centering\arraybackslash}p{1.4cm}
    >{\raggedright\arraybackslash}p{3.4cm}@{}}
\toprule
\textbf{ID} & \textbf{Omschrijving} & \textbf{Score} & \textbf{Klasse} &
\textbf{Kroonjuwelen} \\
\midrule
\endfirsthead
\toprule
\textbf{ID} & \textbf{Omschrijving} & \textbf{Score} & \textbf{Klasse} &
\textbf{Kroonjuwelen} \\
\midrule
\endhead
\midrule
\multicolumn{5}{r}{\emph{vervolgt op volgende pagina}} \\
\endfoot
\bottomrule
\endlastfoot
\textbf{PT1} & Geen rate-limiting op authenticatie-endpoints   & \rood{20} & \rood{Hoog} & KJ3 $\to$ KJ1, KJ2 \\
\addlinespace
\textbf{SQ7} & Ongeauth.\ uitlezing global properties (\texttt{/settings.form/search}) & \rood{20} & \rood{Hoog} & KJ7, KJ4 \\
\addlinespace
\textbf{SQ2} & Onvolledige switch-statements (zoeklogica)      & \rood{16} & \rood{Hoog} & KJ2 \\
\addlinespace
\textbf{SQ1} & Kwetsbare sessie-validatie (session fixation)   & \rood{15} & \rood{Hoog} & KJ3, KJ1, KJ2 \\
\addlinespace
\textbf{SQ3} & Foute testassertions (verbergen bugs)           & \rood{15} & \rood{Hoog} & KJ2 \\
\addlinespace
\textbf{CD1} & Supply chain attack (dependency)                & \rood{15} & \rood{Hoog} & KJ7, KJ1, KJ2 \\
\addlinespace
\textbf{AT2} & Bulk data-exfiltratie via REST API              & \rood{15} & \rood{Hoog} & KJ1, KJ2 \\
\addlinespace
\textbf{SQ8} & Reflected XSS in Swagger debug-endpoint          & \amber{12}$^\triangle$ & \amber{Midden} & KJ3 \\
\addlinespace
\textbf{CD2} & Lekkage van secrets                              & \amber{12}$^\triangle$ & \amber{Midden} & KJ7, KJ3 \\
\addlinespace
\textbf{CQL1}& Log injection (CWE-117) via REST-parameters      & \amber{12} & \amber{Midden} & KJ6 \\
\addlinespace
\textbf{AT3} & Credential-aanvallen via wachtwoordherstel       & \amber{12}$^\triangle$ & \amber{Midden} & KJ3, KJ1, KJ2 \\
\addlinespace
\textbf{AT4} & Gecompromitteerde derde-partij-integratie        & \amber{12}$^\triangle$ & \amber{Midden} & KJ1, KJ2, KJ3 \\
\addlinespace
\textbf{PT2} & Gebruikersenumeratie via `/user` API & \amber{12}$^\triangle$ & \amber{Midden} & KJ3 \\
\addlinespace
\textbf{PT3} & Stack trace disclosure in foutrespons            & \amber{12}$^\triangle$ & \amber{Midden} & KJ6 $\to$ KJ1, KJ2 \\
\end{longtable}
}

\noindent\small $^\triangle$ Midden-score (9--12) maar raakt KJ1/KJ2/KJ3 met
score $\geq 10$ $\to$ conform de risicobereidheid (\S\ref{sec:normenkader})
geclassificeerd als \emph{onacceptabel}.\normalsize

\vspace{0.6em}
Voor twee aanvalsketens is een \textbf{bow-tie-analyse} uitgewerkt
(Bijlage~\ref{app:risicomatrix}): \texttt{SQ1} (sessie-overname, score~15)
en \texttt{CD1} (supply chain, score~15). Twee risico's scoren hoger dan de bow-tie-onderwerpen:
\texttt{SQ7} (ongeauthenticeerde uitlezing van alle global properties, score~20; gemitigeerd via
auth + \texttt{Manage RESTWS}-privilege, commit \texttt{07612ca}) en \texttt{PT1} (geen
rate-limiting, score~20), empirisch bevestigd door de penetratietest
en uitgewerkt als bevinding B-001; de bijbehorende barrière (rate-limiting) maakt bovendien
deel uit van de \texttt{SQ1}-bow-tie, omdat ze zowel sessie-ID- als credential-brute-force
tegenhoudt. De bow-ties \texttt{SQ1} en \texttt{CD1} zijn gekozen als de twee best
\emph{exploiteerbare} aanvalsketens met de breedste impact op de kroonjuwelen; beide
zijn voorzien van preventieve barrières vóór en correctieve maatregelen ná het
top-event.

\section{Bevindingen}
\label{sec:bevindingen}

Hieronder zijn zes bevindingen uitgewerkt, geselecteerd op risicoscore en op
spreiding over de geauditeerde controls. \textbf{Alle zes zijn tijdens de audit
gemitigeerd en hertest}; elke bevinding sluit met een mitigatie- en
validatie-paragraaf met een traceerbare commit-referentie. B-006 (log injection)
toont daarnaast een volledig afgesloten detect-\,$\to$\,mitigeer-\,$\to$\,%
verifieer-cyclus via een CodeQL-herscan. Elke bevinding is handmatig
geverifieerd tegen de broncode (bestand + regel) en, voor de pentestbevindingen,
opnieuw aangevallen op de testomgeving (Kali, 2026-06-16/17;
Bijlage~\ref{app:pentest} \S I.4).

\subsection{B-001 --- Ontbrekende rate-limiting op authenticatie (PT1 / F-01)}
\label{sec:b001}

\begin{tabular}{@{}l >{\raggedright\arraybackslash}p{11.4cm}@{}}
\textbf{Risicoscore} & \rood{20 --- Hoog} (Kans 4 $\times$ Impact 5) --- raakt KJ3 \\
\textbf{NEN-7510}    & A.8.5 Authenticatie \\
\textbf{Kroonjuwelen}& KJ3 --- via overname KJ1, KJ2 \\
\textbf{OWASP}       & A07:2021 Identification \& Authentication Failures (CWE-307) \\
\textbf{Status}      & \groen{Gemitigeerd en hertest --- commit \texttt{59bf421}} \\
\end{tabular}

\textbf{Beschrijving.} De \texttt{AuthorizationFilter} stuurt elke binnenkomende
Basic-Auth-aanvraag rechtstreeks door naar \texttt{Context.authenticate()} zonder
rate-limiting of throttling op netwerkniveau. OpenMRS bevat wel een
\emph{account-lockout} voor bestaande gebruikers (na circa~7 mislukte pogingen), maar
dit geldt \emph{niet} voor niet-bestaande gebruikersnamen en biedt geen bescherming
als het correcte wachtwoord vóór de lockout-drempel wordt gevonden.

\textbf{Bewijs.} \texttt{AuthorizationFilter.java:113}:
\texttt{Context.authenticate(userAndPass[0], userAndPass[1])} wordt in een onbegrensde
stroom aangeroepen; er is geen HTTP-niveau rate-limiter aanwezig. De penetratietest
heeft dit empirisch bevestigd: de geldige admin-credentials (\texttt{admin:Admin123!})
werden gevonden in \textbf{poging~6 van~84}, binnen \textbf{2~seconden}, met een
simpel Python-script, ruim vóór de account-lockout sloeg
(zie Bijlage~\ref{app:pentest}, F-01). Daarnaast werd via timing-analyse bevestigd
dat niet-bestaande accounts nooit geblokkeerd worden (50 opeenvolgende pogingen,
nooit HTTP~429 of vertraging).

\textbf{Impact.} In een ziekenhuisomgeving met veelgebruikte zwakke wachtwoorden
(naam + geboortejaar, afdelingsnaam) is een dictionary attack realistisch en binnen
seconden succesvol. Een gecompromitteerd klinisch account geeft direct toegang tot alle
patiëntdossiers van die gebruiker (KJ1, KJ2). OWASP ASVS Level~1 vereist
rate-limiting als minimumvereiste.

\textbf{Aanbeveling.} Rate-limiting invoeren op de authenticatie-keten:
max.~10~pogingen per minuut per bron-IP, retourneer HTTP~429 bij overschrijding,
gecombineerd met exponentiële back-off. Log overschrijdingen als
\texttt{ACCESS\_BLOCKED} (NEN-7510 A.8.15).

\textbf{Mitigatie.} Een in-memory IP-rate-limiter is toegevoegd aan
\texttt{AuthorizationFilter}: een begrensde \texttt{ConcurrentHashMap} houdt per
bron-IP het aantal mislukte pogingen binnen een venster van 60~seconden bij. Bij
meer dan 10 mislukte pogingen retourneert het filter \texttt{HTTP 429} mét een
\texttt{return} vóór \texttt{Context.authenticate()}, en logt het
\texttt{ACCESS\_BLOCKED rate\_limited=true} op WARN; de teller reset bij een
succesvolle authenticatie. Geen extra dependency (Java 8-compatibel). Commit
\texttt{59bf421}.

\textbf{Validatie.} Pentest-hertest op Kali (2026-06-16): de 11e opeenvolgende
mislukte Basic-Auth-poging gaf \texttt{HTTP 429} (vóór de fix: nooit geblokkeerd),
met de bijbehorende \texttt{ACCESS\_BLOCKED}-logregel. Zes nieuwe unit tests
(\texttt{AuthorizationFilterRateLimitTest}) en de twaalf bestaande
filtertests slagen; \texttt{mvn clean install} groen. Zie
\texttt{06-mitigatie-en-validatie/06.md} \S6 en Bijlage~\ref{app:pentest} \S I.4.

\subsection{B-002 --- Sessie-fixatie en best-effort toegangsfilter (SQ1 / F-03)}
\label{sec:b002}

\begin{tabular}{@{}l >{\raggedright\arraybackslash}p{11.4cm}@{}}
\textbf{Risicoscore} & \rood{15 --- Hoog} (Kans 3 $\times$ Impact 5) \\
\textbf{NEN-7510}    & A.8.5 Authenticatie (tevens A.8.3 Toegangsbeveiliging) \\
\textbf{Kroonjuwelen}& KJ3 --- via overname KJ1, KJ2 \\
\textbf{OWASP}       & A07:2021 Identification \& Authentication Failures (CWE-384/302) \\
\textbf{Status}      & \groen{Gemitigeerd en hertest} --- CWE-384 (sessie-rotatie) commit \texttt{9c2fb64}, CWE-302 (fall-through) commit \texttt{e3f3c7d} \\
\end{tabular}

\textbf{Beschrijving (vóór fix).} De \texttt{AuthorizationFilter} ververste het
sessie-ID \emph{niet} na een succesvolle login: een door een aanvaller vooraf
``geplant'' sessie-ID bleef na authenticatie geldig (session fixation, CWE-384).
De filter is bovendien \emph{by design} best-effort --- het klassecommentaar stelt
dat het filter \emph{``never stops execution''} en op de service-laag vertrouwt om
alsnog te falen.

\textbf{Bewijs (vóór fix).} \texttt{AuthorizationFilter.java:82}:
\texttt{getRequestedSessionId()} gaf het client-aangeleverde sessie-ID terug
zonder server-side validatie; na \texttt{Context.authenticate()} werd de sessie
niet geïnvalideerd (CWE-384). Daarnaast liep een afzonderlijke CWE-302
fall-through: de ontbrekende \texttt{return} na de \texttt{sendError(401)} bij een
verlopen sessie liet de keten dóórlopen naar \texttt{chain.doFilter()}.

\textbf{Impact.} Sessie-overname geeft een aanvaller volledige toegang tot
patiëntdossiers (KJ1, KJ2) met de rechten van het slachtoffer; dit is een meldplichtig
datalek (AVG Art.~33). In combinatie met B-004 (XSS) ontstaat een complete keten:
XSS steelt de sessiecookie, waarna de niet-verversde sessie wordt overgenomen.

\textbf{Mitigatie.} Twee maatregelen, beide in \texttt{AuthorizationFilter}:
(1)~na een succesvolle \texttt{Context.authenticate()} wordt de bestaande sessie
geïnvalideerd en een nieuwe aangemaakt (\texttt{getSession(false)} $\to$
\texttt{invalidate()} $\to$ \texttt{getSession(true)}), zodat Tomcat een nieuw
\texttt{JSESSIONID} uitgeeft (CWE-384, commit \texttt{9c2fb64}); (2)~een
\texttt{return} is toegevoegd direct na de \texttt{sendError(401)} bij sessie-time-out,
zodat de keten daadwerkelijk stopt (CWE-302, commit \texttt{e3f3c7d}). De
CodeQL-melding \texttt{java/user-controlled-bypass} op de optionele
Basic-Auth-tak (alert \#908/\#913) staat na de wijziging op \texttt{state: fixed}
(geaccepteerd risico / false positive; analyse in
\texttt{06-mitigatie-en-validatie/06.md} \S3b).

\textbf{Validatie.} Pentest-hertest op Kali (2026-06-17): het \texttt{JSESSIONID}
verschilde nu vóór en ná de login (vóór de fix: identiek). Vijf nieuwe tests
(\texttt{AuthorizationFilterSessionFixationTest}) en de bestaande twaalf
filtertests slagen; \texttt{mvn clean install} groen. Op termijn blijft
declaratieve autorisatie (Spring Security) ter vervanging van het best-effort
filter aanbevolen (roadmap). Zie \texttt{06.md} \S13 en Bijlage~\ref{app:pentest} \S I.4.

\subsection{B-003 --- Gebruikersenumeratie via `/user` API (PT2 / F-02)}
\label{sec:b003}

\begin{tabular}{@{}l >{\raggedright\arraybackslash}p{11.4cm}@{}}
\textbf{Risicoscore} & \amber{12 --- Midden}$^\triangle$ (Kans 4 $\times$ Impact 3) --- onacceptabel (raakt KJ3 met score $\geq 10$) \\
\textbf{NEN-7510}    & A.8.5 Authenticatie \\
\textbf{Kroonjuwelen}& KJ3 --- via misbruik indirect KJ1, KJ2 \\
\textbf{OWASP}       & A07:2021 Identification \& Authentication Failures (CWE-204) \\
\textbf{Status}      & \groen{Gemitigeerd en hertest --- commit \texttt{ed6f3da}} \\
\end{tabular}

\textbf{Beschrijving.} Het endpoint \texttt{GET /openmrs/ws/rest/v1/user}
retourneert bij geauthenticeerde toegang een volledige lijst van alle
gebruikersaccounts in het systeem, inclusief gebruikersnamen en persoonlijke
details. Er is geen fijnmazige autorisatiecontrole: zelfs gebruikers met minimale
privileges (bv.\ laboratorium-technicus) kunnen alle systeembeheerdersaccounts
enumereren.

\textbf{Bewijs.} Pentest F-02 (11--12 juni 2026):
\begin{lstlisting}[basicstyle=\ttfamily\scriptsize, breaklines=true, frame=single]
GET /openmrs/ws/rest/v1/user
Authorization: Basic <laag-geprivilegieerde-creds>

--> Retourneert alle gebruikers:
{
  "results": [
    { "display": "admin", "username": "admin", ... },
    { "display": "doctor", "username": "doctor", ... },
    ...
  ]
}
\end{lstlisting}

\textbf{Impact.} Een aanvaller verkrijgt een complete lijst van geldige
gebruikersnamen en systeemdetails. Dit maakt gerichte brute-force aanvallen (B-001, F-01)
significant efficiënter: in plaats van willekeurige namen te proberen kan een
aanvaller direct de verkregen accounts targeten. In combinatie zijn B-001 en B-003
een complete credential-compromis-keten.

\textbf{Aanbeveling.} Beperk de \texttt{/user}-API tot expliciete beheerdersrollen;
controleer fijnmazige autorisatie per resource; retourneer geen gevoelige
accountdetails aan lager-geprivilegieerde gebruikers.

\textbf{Mitigatie.} Twee maatregelen: (A)~\texttt{Context.requirePrivilege(%
PrivilegeConstants.GET\_USERS)} als eerste instructie van \texttt{doGetAll()} en
\texttt{doSearch()} in \texttt{UserResource1\_8}, zodat zowel het lijsten als het
zoeken op \texttt{/user} nu het OpenMRS-privilege \emph{Get Users} vereist; (B)~de
\texttt{AUTH\_FAILURE}-logregel in \texttt{AuthorizationFilter} logt uniform,
zonder onderscheid meer tussen fout wachtwoord en account-lockout (sluit een
side-channel). Commit \texttt{ed6f3da}.

\textbf{Validatie.} Pentest-hertest op Kali (2026-06-16): \texttt{GET /user} en
\texttt{GET /user?q=...} als laag-geprivilegieerde gebruiker worden nu geweigerd
(\emph{``Privileges required: Get Users''}) i.p.v.\ de volledige lijst te
retourneren; de auth-logregels zijn identiek ongeacht of het account bestaat. Drie
nieuwe unit tests slagen; \texttt{mvn clean install} groen. Zie
\texttt{06.md} \S7 en Bijlage~\ref{app:pentest} \S I.4.

\subsection{B-004 --- Reflected XSS in het Swagger debug-endpoint (SQ8): gemitigeerd}
\label{sec:b004}

\begin{tabular}{@{}l >{\raggedright\arraybackslash}p{11.4cm}@{}}
\textbf{Risicoscore} & \amber{12 --- Midden}$^\triangle$ (Kans 3 $\times$ Impact 4) --- vóór mitigatie \\
\textbf{NEN-7510}    & A.8.28 Veilig coderen (output-encoding), A.8.26 \\
\textbf{Kroonjuwelen}& KJ3 --- via overname KJ1, KJ2 \\
\textbf{OWASP}       & A03:2021 Injection (XSS) --- CodeQL \texttt{java/xss}, alert \#1 (high) \\
\textbf{Status}      & \groen{Gemitigeerd --- commit \texttt{e3f3c7d}} \\
\end{tabular}

\textbf{Beschrijving (vóór fix).} De \texttt{debug}-handler plaatste de
request-parameter \texttt{tag} ongesanitized in een HTML-respons:
\begin{quote}\ttfamily\small return "<h1>Debugging Tag: " + tag + "</h1>";\end{quote}
Een aanvaller kon \texttt{?tag=<script>\ldots</script>} injecteren. Het endpoint
\texttt{/module/\allowbreak webservices/\allowbreak rest/\allowbreak apiDocs/\allowbreak debug}
ligt buiten de \texttt{AuthorizationFilter}-keten en was dus zonder authenticatie
én zonder IP-allowlist bereikbaar.

\textbf{Bewijs (vóór fix).} \texttt{SwaggerDocController.java:27} (CodeQL-meldtekst:
\emph{``Cross-site scripting vulnerability due to a user-provided value''}).

\textbf{Impact (vóór fix).} Eén crafted phishing-link volstaat: uitgevoerd in de
browser van een ingelogde admin/zorgverlener kon de XSS sessiecookies of
CSRF-tokens stelen $\to$ overname van een geprivilegieerde sessie (KJ3) en van
daaruit toegang tot KJ1/KJ2. Geen voorafgaande toegang vereist.

\textbf{Mitigatie (twee lagen).} \texttt{tag} wordt nu HTML-geëncodeerd met
\texttt{org.springframework.web.util.HtmlUtils.htmlEscape(...)} vóórdat het in
de respons wordt geplaatst; dit neutraliseert de \texttt{java/xss}-taintflow.
Als defense-in-depth is de hele \texttt{/apiDocs}-boom bovendien uitschakelbaar
via de global property \texttt{webservices.rest.enableSwaggerDocs}: in productie
(\texttt{false}) geeft het endpoint HTTP 404 en is de vector volledig weg.

\textbf{Validatie (statisch + unit test).} SQ8 is --- anders dan de
pentestbevindingen B-001--B-003/B-005 --- \emph{niet via een dynamische pentest}
gevalideerd: het \texttt{/apiDocs/debug}-endpoint staat in de testomgeving uit
(\texttt{enableSwaggerDocs=false}) en geeft HTTP~404, waardoor een live
curl-/browser-hertest niet mogelijk is. De validatie steunt op twee statische,
reproduceerbare bronnen. \emph{(1)}~De automatische CodeQL-herscan op \texttt{main}
(run \texttt{27471405779}, commit \texttt{b5627faa} via PR~\#96, 2026-06-13) zet
\texttt{java/xss} alert~\#1 op \texttt{state: fixed}, omdat \texttt{HtmlUtils.htmlEscape}
de door CodeQL herkende taint-barrier is. \emph{(2)}~Een dedicated unit test
\texttt{SwaggerDocControllerTest} roept \texttt{debug()} rechtstreeks aan (zonder
HTTP-laag, dus onafhankelijk van de endpoint-404) en bewijst deterministisch dat
\texttt{<script>alert(1)</script>} als \texttt{\&lt;script\&gt;} wordt geëncodeerd én dat
de gating bij uitgeschakelde docs HTTP~404 zonder body geeft --- \textbf{2 tests groen}
(\texttt{Tests run: 2, Failures: 0, Errors: 0}). \texttt{mvn clean install} groen, geen
regressie. Resterende beperking: geen dynamische pentest-hertest (endpoint 404 in de
testomgeving).

\subsection{B-005 --- Stack trace disclosure in foutrespons (PT3 / F-04)}
\label{sec:b005}

\begin{tabular}{@{}l >{\raggedright\arraybackslash}p{11.4cm}@{}}
\textbf{Risicoscore} & \amber{12 --- Midden}$^\triangle$ (Kans 4 $\times$ Impact 3) --- onacceptabel (enabler voor CVE-targeting) \\
\textbf{NEN-7510}    & A.8.15 Logging (tevens A.8.8 Kwetsbaarheden) \\
\textbf{Kroonjuwelen}& KJ6 --- indirect KJ1, KJ2 (reconnaissance-enabler) \\
\textbf{OWASP}       & A09:2025 Security Logging and Alerting Failures (CWE-209) \\
\textbf{Status}      & \groen{Gemitigeerd en hertest --- commit \texttt{9c30cdd}} \\
\end{tabular}

\textbf{Beschrijving.} Bij unauthenticated toegang en interne fouten retourneert
OpenMRS REST volledige Java stack traces in de HTTP response body. Drie endpoints
zijn bevestigd kwetsbaar: \texttt{GET /user}, \texttt{GET /module} en
\texttt{POST /user}. De stack traces bevatten achtereenvolgens: de exacte
autorisatielaag met regelnummer (\texttt{org.openmrs.aop.AuthorizationAdvice:120}),
de Hibernate-versie en ORM-configuratiestructuur, de Jackson-dependencyversie en
het interne domeinmodel (\texttt{org.openmrs.Role.getId(Role.java:246)}).

\textbf{Bewijs.} Pentest F-04 (11--12 juni 2026); bevestigd met Burp Suite en
\texttt{curl}:
\begin{lstlisting}[basicstyle=\ttfamily\scriptsize, breaklines=true, frame=single]
curl -s 'http://192.168.56.1:8080/openmrs/ws/rest/v1/user'
curl -s 'http://192.168.56.1:8080/openmrs/ws/rest/v1/module'
curl -u 'admin:Admin123!' -X POST 'http://192.168.56.1:8080/openmrs/ws/rest/v1/user'
# Alle drie retourneren ~80 regels Java stack trace in de response body
\end{lstlisting}
Screenshots in Bijlage~\ref{app:pentest} (Figuren 4--6).

\textbf{Impact.} Stack traces maken vervolgaanvallen significant eenvoudiger:
exacte framework-versies geven een aanvaller directe CVE-targeting (synergie met
SBOM-bevindingen in Hoofdstuk~\ref{ch:sbom}), de ORM-configuratie verlaagt de
drempel voor deserialisatie-analyse en de autorisatiearchitectuur versnelt
privilege-escalation-planning. De bevinding is laagdrempelig te reproduceren
zonder credentials.

\textbf{Aanbeveling.} Configureer een globale \texttt{@ExceptionHandler} die
interne details vervangt door een generiek foutbericht met een referentie-ID.
Log de volledige stack trace uitsluitend intern (NEN-7510 A.8.15); stuur
nooit interne details naar de HTTP-client.

\textbf{Mitigatie.} Twee maatregelen: (A)~in \texttt{RestUtil.wrapErrorResponse()}
zijn het \texttt{code}-veld (voorheen \texttt{className:lineNumber}) en het
\texttt{detail}-veld (voorheen de volledige stack trace) nu altijd leeg; de global
property \texttt{webservices.rest.enableStackTraceDetails} wordt bewust genegeerd.
(B)~In \texttt{BaseRestController} wordt \texttt{ContextAuthenticationException}
nu vóór \texttt{APIAuthenticationException} herkend en als \texttt{HTTP 403/401}
afgehandeld; overige 5xx-fouten krijgen een UUID-correlatie-ID dat intern wordt
gelogd en als \texttt{code} wordt teruggegeven, met een leeg \texttt{detail}-veld.
Commit \texttt{9c30cdd}.

\textbf{Validatie.} Pentest-hertest op Kali (2026-06-16): \texttt{GET /user} als
laag-geprivilegieerde gebruiker gaf een leeg \texttt{detail}-veld en \texttt{HTTP
403} (vóór de fix: \texttt{className:lineNumber} + $\sim$60 regels stack trace +
\texttt{HTTP 500}); met \texttt{enableStackTraceDetails=true} blijft \texttt{detail}
leeg. Bijgewerkte en nieuwe tests in \texttt{BaseRestControllerTest} en
\texttt{RestUtilTest} slagen; \texttt{mvn clean install} groen. Zie
\texttt{06.md} \S8 en Bijlage~\ref{app:pentest} \S I.4.

\subsection{B-006 --- Log injection (CWE-117): gemitigeerd en geverifieerd (CQL1 / F-07)}
\label{sec:b006}

\begin{tabular}{@{}l >{\raggedright\arraybackslash}p{11.4cm}@{}}
\textbf{Risicoscore} & \amber{12 --- Midden} (Kans 4 $\times$ Impact 3) --- raakt KJ6 \\
\textbf{NEN-7510}    & A.8.15 Logging \\
\textbf{Kroonjuwelen}& KJ6 Audit logs (integriteit) \\
\textbf{OWASP}       & A09:2021 Logging \& Monitoring Failures (CWE-117) \\
\textbf{Status}      & \groen{Gemitigeerd en geverifieerd --- commit \texttt{831dafb}} \\
\end{tabular}

\textbf{Beschrijving.} Meerdere REST-controllers logden request-afgeleide waarden
(\texttt{resource}, \texttt{uuid}, query-parameters) zonder neutralisatie. Een
aanvaller kon CR/LF-tekens injecteren om valse logregels te vervalsen (``log
forging''), bijvoorbeeld een nep-\texttt{INFO}-regel die een geslaagde login
suggereert --- waarmee forensisch onderzoek wordt misleid en de integriteit van het
audit trail (KJ6) wordt aangetast.

\textbf{Bewijs (vóór fix).} CodeQL \texttt{java/log-injection}
(\texttt{security-extended}): 9~alerts in o.a.\ \texttt{RestUtil.java:521},
\texttt{ClearDbCacheController2\_0.java:79/86/95} en
\texttt{ConceptResource1\_8.java:567}.

\textbf{Mitigatie.} Eén centrale helper \texttt{RestUtil.sanitizeForLog(String)}
verwijdert \texttt{\textbackslash r}/\texttt{\textbackslash n} uit user-controlled
waarden (het door CodeQL herkende taint-barrier-patroon) en is toegepast op alle
acht call-sites. Gerealiseerd met Claude Code (Sonnet~4.6); de fix is handmatig
gereviewd vóór commit.

\textbf{Validatie (voor/na).} De CodeQL-run op \texttt{main}
(commit \texttt{e476489}, 2026-06-10) sluit \textbf{9 van 9}
\texttt{java/log-injection}-alerts (\#7--15, \texttt{state: fixed}). Open alerts:
\textbf{9 $\to$ 0}. Daarmee is de detect-\,$\to$\,mitigeer-\,$\to$\,verifieer-cyclus
kwantitatief en reproduceerbaar afgesloten.

\textbf{Vervolg (gesloten).} Dezelfde herschrijving introduceerde 4 \emph{nieuwe}
\texttt{java/log-injection}-alerts (\texttt{AuthorizationFilter.java:72/84/115/122},
alert-nrs \#808--811), aanvankelijk als \emph{``false positive''} gedismissed
\emph{zonder} onderbouwing. Omdat die classificatie afhing van de (in deze module
afwezige) log-appender-config en dus niet aantoonbaar was, zijn de user-controlled
waarden (request-URI, \texttt{attemptedUser}, exception-message) alsnog via
\texttt{RestUtil.sanitizeForLog()} gesanitized (commit \texttt{07612ca}). De vier
alerts staan na de CodeQL-scan van \texttt{f780f28} (2026-06-16) op
\texttt{state: fixed}; daarmee zijn alle 13 \texttt{java/log-injection}-alerts
(de oorspronkelijke 9 + deze 4) gesloten.

\newpage

% ════════════════════════════════════════════════════════════════════════════
%  5. SBOM & SUPPLY CHAIN SECURITY
% ════════════════════════════════════════════════════════════════════════════
\chapter{SBOM en Supply Chain Security}
\label{ch:sbom}

Een aanzienlijk deel van het aanvalsoppervlak van de module zit niet in de eigen
broncode, maar in de \emph{toeleveringsketen}: 237 externe open-source componenten
die worden meegebundeld in het productie-artefact. Dit hoofdstuk inventariseert
die keten (\S\ref{sec:sbom}), vat de belangrijkste bekende kwetsbaarheden samen
(\S\ref{sec:top-dependencies}) en licht toe waarom de prioritering op een
\emph{contextuele} risicoscore berust in plaats van op de ruwe CVSS-score
(\S\ref{sec:contextuele-risico}). De volledige analyse --- inclusief de
geprioriteerde backlog van 16 bevindingen en het patch-/suppress-/accept-besluit
per kwetsbaarheid --- is vastgelegd in het bronbestand \texttt{03.md} onder
\texttt{docs/security/03-sbom-en-cve-advies/}.

\section{Software Bill of Materials}
\label{sec:sbom}

De SBOM is het fundament van kwetsbaarhedenbeheer (NEN-7510:2024-2 \textbf{A.8.8}):
zonder een actuele, machine-leesbare inventaris van alle componenten is niet vast
te stellen wélke onderdelen gepatcht moeten worden. De module genereert deze
inventaris geautomatiseerd in het CI-proces.

\begin{tabular}{@{}ll}
    \toprule
    \textbf{Eigenschap} & \textbf{Waarde} \\
    \midrule
    Formaat          & CycloneDX 1.6 (JSON) \\
    Artefact         & \texttt{sbom.cdx.json} ($\sim$520 kB) \\
    Componenten      & 238 (237 afhankelijkheden + de module zelf) \\
    Generator        & \texttt{cyclonedx-maven-plugin:2.9.1:makeAggregateBom} \\
    Trigger          & elke push naar \texttt{main}/\texttt{develop} + elke PR \\
    CI-workflow      & \texttt{.github/workflows/sbom.yml} \\
    CI-artefact      & \texttt{sbom-cyclonedx} (90 dagen bewaard) \\
    Aanvullende scan & OSV-Scanner + Snyk in dezelfde job \\
    \bottomrule
\end{tabular}

De plugin-versie is \emph{gepind} (\texttt{2.9.1}) zodat de SBOM-generator zelf
reproduceerbaar is en niet stilzwijgend naar een nieuwere, niet-vertrouwde release
wordt opgewaardeerd. CycloneDX is bovendien het formaat dat de \textbf{Cyber
Resilience Act (CRA)} als (toekomstige) wettelijke eis aan softwareproducten
hanteert; de module voldoet daarmee nu al aan die aankomende verplichting. Het
bestand zelf is opgenomen als Bijlage~\ref{app:sbom}.

\subsection*{Reikwijdte van de SBOM}

De \texttt{makeAggregateBom}-doelstelling neemt standaard alléén de \texttt{compile}-,
\texttt{runtime}- en \texttt{provided}-scopes mee; de \texttt{test}-scope wordt
uitgesloten. De 238 componenten vormen daarmee de \emph{productie-relevante}
afhankelijkheidsset. Dit is geverifieerd: test-only componenten zoals de
H2-testdatabase en JUnit/Mockito komen niet in de SBOM voor. Dat is op zichzelf
bewijs voor het suppress-besluit rond de twee H2-RCE's
(\S\ref{sec:contextuele-risico}): de component bereikt het productie-artefact
aantoonbaar niet.

\subsection*{Licentie-observatie}

Van de 237 afhankelijkheden staan er 160 onder Apache-2.0 en nog enkele tientallen
onder permissieve licenties (BSD, MIT, UPL). Daarnaast bevat de keten echter ook
\emph{copyleft}-componenten --- LGPL-2.1, GPL en MPL-1.1 --- en negen componenten
zonder gedeclareerde licentie. Voor een zorgproduct dat (her)gedistribueerd kan
worden is dat een aandachtspunt voor de juridische supply-chain-naleving, naast
het puur technische kwetsbaarhedenbeheer.

\section{Top Dependencies en Bekende CVE's}
\label{sec:top-dependencies}

Drie onafhankelijke SCA-bronnen zijn gekruisvalideerd: \textbf{Dependabot}
(GitHub Advisory Database), \textbf{OSV-Scanner} (scant de CycloneDX SBOM) en
\textbf{Snyk} (eigen database). De \textbf{264 ruwe Dependabot-alerts} herleiden
zich na deduplicatie tot \textbf{75 unieke CVE's} (3 critical, 32 high, 35 medium,
5 low); OSV-Scanner bevestigt daarvan 64 onafhankelijk en Snyk voegt nog 36 unieke
CVE's toe die Dependabot mist. De onderstaande tabel toont de hoogst scorende
bevindingen ná contextuele weging (\S\ref{sec:contextuele-risico}); de volledige
top-16-backlog staat in \texttt{03.md} \S4 en het Snyk-rapport in
Bijlage~\ref{app:snyk}.

{\small
\begin{longtable}{@{}
    >{\raggedright\arraybackslash}p{3.6cm}
    >{\raggedright\arraybackslash}p{1.2cm}
    >{\centering\arraybackslash}p{0.85cm}
    >{\centering\arraybackslash}p{0.95cm}
    >{\raggedright\arraybackslash}p{1.4cm}
    >{\raggedright\arraybackslash}p{5.4cm}@{}}
\toprule
\textbf{Component} & \textbf{Versie} & \textbf{CVSS} & \textbf{Score} &
\textbf{Besluit} & \textbf{Toelichting} \\
\midrule
\endfirsthead
\toprule
\textbf{Component} & \textbf{Versie} & \textbf{CVSS} & \textbf{Score} &
\textbf{Besluit} & \textbf{Toelichting} \\
\midrule
\endhead
\midrule
\multicolumn{6}{r}{\emph{vervolgt op volgende pagina}} \\
\endfoot
\bottomrule
\endlastfoot
\texttt{spring-web} & 5.3.30 & 9.8 & \rood{9.0} & Patch &
    CVE-2016-1000027; deserialisatie\,$\rightarrow$\,RCE, kern van de REST-laag.
    Fix = Spring~6 (major, gefaseerd). \\
\addlinespace
\texttt{jackson-mapper-asl} & 1.9.14 & 9.8 & \rood{8.1} & Patch &
    Codehaus Jackson is EOL; migreren naar \texttt{jackson-databind}~2.x. \\
\addlinespace
\texttt{snakeyaml} & 1.30 & 8.3 & \rood{6.2} & Patch &
    CVE-2022-1471, hoogste EPSS (0.94). Zie convergentie-observatie hieronder. \\
\addlinespace
\texttt{hibernate-core} & 5.6.15 & 8.3 & \amber{6.0} & Accept &
    CVE-2026-0603 (SQLi); geen patch, EPSS verwaarloosbaar, parametrized criteria. \\
\addlinespace
\texttt{netty-handler} & 4.1.118 & 9.2 & \amber{5.9} & Patch &
    HTTP-I/O-laag; fix in 4.1.119+ (non-breaking). \\
\addlinespace
\texttt{c3p0} & 0.9.5.5 & 8.9 & \amber{5.8} & Patch &
    Deserialisatie in de DB-connection-pool; fix in 0.10.x. \\
\addlinespace
\texttt{spring-web} & 5.3.30 & 8.1 & \amber{5.7} & Patch &
    CVE-2024-22262; patch 5.3.34 \emph{zonder} breaking changes. \\
\addlinespace
\texttt{log4j-core} & 2.22.1 & 7.7 & \amber{4.7} & Patch &
    Log-injectie (CWE-117); versterkt de applicatie-fix uit commit \texttt{831dafb}. \\
\end{longtable}
}

\subsection*{SBOM-observatie: gemengde \texttt{snakeyaml}-versies}

De SBOM legt een concreet supply-chain-defect bloot dat een losse CVE-lijst zou
missen: \texttt{snakeyaml} komt \emph{tweemaal} in de afhankelijkheidsboom voor ---
zowel de kwetsbare \texttt{1.30} als de reeds gepatchte \texttt{2.4}, via
verschillende transitieve paden. Het patchen van bevinding~\#3 is daarmee geen
simpele versie-bump maar een \emph{convergentieprobleem}: pas een geforceerde
\texttt{<dependencyManagement>}-override naar één 2.x-versie verwijdert het
kwetsbare \texttt{1.30}-artefact volledig uit de productie-build.

\section{Contextuele Risicoanalyse}
\label{sec:contextuele-risico}

Een ruwe CVSS-score zegt niets over het \emph{werkelijke} risico voor déze module.
Daarom is elke bevinding gescoord met een contextuele formule die CVSS combineert
met drie genormaliseerde factoren (schaal 0.0--1.0):

\begin{quote}
\textbf{Contextueel = CVSS $\times$ (Bereikbaarheid $\times$ 0.4 +
Healthcare-impact $\times$ 0.4 + Exploiteerbaarheid $\times$ 0.2)}
\end{quote}

\textbf{Bereikbaarheid} weegt of de component in het REST-aanvraagpad zit
(runtime/productie) of alleen in test/build; \textbf{healthcare-impact} weegt of
misbruik patiënt-/medische gegevens of authenticatie raakt; \textbf{exploiteerbaarheid}
gebruikt de EPSS-waarde (kans op misbruik binnen 30 dagen). Geen van de 75 CVE's
staat in de CISA KEV-catalogus (gecontroleerd 2026-06-08) --- er is dus geen
\emph{actief} misbruik in het wild bekend voor deze versies. De volledige
scoringscriteria per factor (de drempelwaarden voor elke stap op de
0.0--1.0-schaal), de berekende score per bevinding en het bijbehorende
patch-/suppress-/accept-besluit zijn uitgewerkt in het bronbestand \texttt{03.md}
\S3--5.

Het effect van deze weging is fundamenteel. De twee H2-database-RCE's hebben de
\emph{hoogst mogelijke} CVSS (9.8) én een zeer hoge EPSS (0.91 / 0.27), maar zakken
naar de onderkant van de backlog: H2 draait uitsluitend in \texttt{test}-scope en
komt --- zoals in \S\ref{sec:sbom} aangetoond --- niet eens in de productie-SBOM
voor. Omgekeerd stijgt \texttt{spring-web} (CVE-2016-1000027), met een
\emph{gemiddelde} EPSS van 0.60, naar de absolute top: het is zowel maximaal
bereikbaar (kern van de REST-laag) als maximaal impactvol (deserialisatie
$\rightarrow$ RCE $\rightarrow$ volledige toegang tot de kroonjuwelen KJ1/KJ2/KJ3).

Dit is de kern van de audit-mindset: niet de scanner-output overnemen, maar
interpreteren. Het advies prioriteert dan ook \textbf{vijf non-breaking quick-wins}
(\texttt{snakeyaml}, \texttt{spring-web}~5.3.34, \texttt{netty-handler},
\texttt{c3p0}, \texttt{log4j-core}) boven de twee grote EOL-migraties
(\texttt{spring-web}~6 en \texttt{struts-core}-vervanging), die --- met onderbouwing
--- als geaccepteerde kortetermijnrisico's met interim-mitigatie worden vastgelegd.

\newpage

% ════════════════════════════════════════════════════════════════════════════
%  6. CONCLUSIE & ADVIES
% ════════════════════════════════════════════════════════════════════════════
\chapter{Conclusie en Advies}
\label{ch:conclusie}

\section{Antwoord op de Auditvraag}
\label{sec:auditvraag}

De auditvraag was of de OpenMRS REST Webservices Module (v3.2.0) voldoet aan
de relevante NEN-7510:2024-2 controls en aan de CRA-verplichtingen voor
open-source componenten in de toeleveringsketen.

\textbf{NEN-7510:2024-2.} Van de zeven controls uit het Normenkader
(\S\ref{sec:normenkader}) voldoet de module na de tijdens deze audit
doorgevoerde mitigaties aan vijf (A.8.5, A.8.8, A.8.9, A.8.15, A.8.25) en is
er bij twee sprake van gedeeltelijke compliance (A.8.3 en A.8.24) --- zie
\S\ref{sec:overall-status} voor het overzicht en
Bijlage~\ref{app:traceability} voor de onderliggende bewijsvoering. A.8.5
(authenticatie) en A.8.9 (configuratiebeheer) zijn van \emph{tekort} naar
\emph{geïmplementeerd} gebracht doordat de onderliggende bevindingen B-001
(rate-limiting, score 20), B-002 (sessie-fixatie) en de ongeauthenticeerde
configuratie-uitlezing (settings.form, score 20) zijn gemitigeerd en hertest.
A.8.24 (cryptografie/TLS) blijft gedeeltelijk: TLS-afdwinging hoort op
deployment-niveau en valt buiten de module-grens. Daarnaast zijn tijdens deze
audit zeven \emph{aanvullende} controls geïmplementeerd en getest
(least-privilege database, Swagger/Host-validatie, supply-chain-signing,
secure coding, OTAP-scheiding en beveiligingstests) die niet in het
oorspronkelijke Normenkader stonden maar wel uit de risicoanalyse volgden.

\textbf{CRA.} De CycloneDX-SBOM (238 componenten, Bijlage~\ref{app:sbom}) en
de gekruisvalideerde SCA-keten (Dependabot + OSV-Scanner + Snyk, 75 unieke
CVE's, \S\ref{sec:top-dependencies}) geven de toeleveringsketen een
machine-leesbare, traceerbare inventarisatie --- een kernvereiste van de CRA.
De contextuele risicoanalyse (\S\ref{sec:contextuele-risico}) onderbouwt per
kwetsbaarheid een patch-/suppress-/accept-besluit; twee non-breaking
quick-win-patches (\texttt{snakeyaml}, \texttt{spring-web}) staan nog open.
De nieuwe signed-release-workflow (cosign keyless-signing + SLSA-provenance,
Bijlage~\ref{app:traceability} A.5.23/A.8.32) levert daarnaast verifieerbare
build-provenance, wat de CRA-conforme transparantie van de supply chain
verder versterkt.

\textbf{Samenvattend:} de module is na deze audit \amber{grotendeels, maar nog
niet volledig compliant} met NEN-7510:2024-2. Alle zes uitgewerkte
code-bevindingen (B-001 t/m B-006) plus de ongeauthenticeerde
configuratie-uitlezing zijn gemitigeerd en hertest, waarmee de eerdere
tekortkomingen op A.8.5 en A.8.9 zijn weggewerkt. De resterende afstand zit
niet in de modulecode maar in (i)~de toeleveringsketen --- twee
non-breaking quick-win-updates (\texttt{snakeyaml}, \texttt{spring-web}) en
enkele geplande EOL-migraties staan nog open --- en (ii)~A.8.24, de
TLS-afdwinging op deployment-niveau. De CRA-gerelateerde eisen aan
supply-chain-transparantie worden grotendeels ondersteund door de
SBOM/SCA/signing-keten die tijdens deze audit is opgezet en geverifieerd. Met
de in \S\ref{sec:roadmap} geprioriteerde roadmap is volledige compliance op de
geauditeerde controls binnen 1--2 sprints realistisch.

\section{Geprioriteerde Aanbevelingen}
\label{sec:aanbevelingen}

De aanbevelingen zijn geordend op urgentie conform de risicobereidheid
(\S\ref{sec:risicomatrix}). \emph{Nu} = score $\geq 15$ of onacceptabel,
directe actie; \emph{Deze sprint} = score 9--12, mitigatie binnen 1~maand;
\emph{Later} = geplande roadmap.

\vspace{0.4em}
{\small
\begin{tabularx}{\textwidth}{@{}lXll@{}}
\toprule
\textbf{Tijdlijn} & \textbf{Actie} & \textbf{Bev.} & \textbf{Control} \\
\midrule
\groen{Gemitigeerd} &
    Rate limiting in \texttt{AuthorizationFilter}: max.~10 mislukte
    pogingen/min per bron-IP, HTTP~429, \texttt{ACCESS\_BLOCKED}-logging
    (commit \texttt{59bf421}) &
    B-001 & A.8.5 \\
\groen{Gemitigeerd} &
    Sessie-ID geregenereerd na login (\texttt{invalidate()} + nieuwe sessie)
    en \texttt{return} na 401 bij time-out (commits \texttt{9c2fb64},
    \texttt{e3f3c7d}) &
    B-002 & A.8.5 \\
\groen{Gemitigeerd} &
    \texttt{/user}-API achter \emph{Get Users}-privilege + uniforme
    \texttt{AUTH\_FAILURE}-logging (commit \texttt{ed6f3da}) &
    B-003 & A.8.5 \\
\groen{Gemitigeerd} &
    Stack traces intern gelogd, generiek bericht + referentie-ID naar de
    client; \texttt{ContextAuthenticationException} $\to$ HTTP 403
    (commit \texttt{9c30cdd}) &
    B-005 & A.8.15 \\
\groen{Gemitigeerd} &
    Ongeauthenticeerde \texttt{settings.form/search} afgeschermd met
    authenticatie + \emph{Manage RESTWS}-privilege (commit \texttt{07612ca}) &
    --- & A.8.9 \\
\rood{Nu} &
    Quick-win component-updates: \texttt{snakeyaml} $\to$ 2.x,
    \texttt{spring-web} $\to$ 5.3.34 &
    SBOM & A.8.8 \\
\amber{Deze sprint} &
    TLS-afdwinging op deployment-niveau (reverse proxy/servlet-container)
    documenteren en verifiëren &
    --- & A.8.24 \\
\geel{Later} &
    Declaratieve autorisatie (Spring Security) ter vervanging van
    het best-effort \texttt{AuthorizationFilter} &
    B-002 & A.8.5 \\
\geel{Later} &
    EOL-migraties: \texttt{spring-web}~6, Tomcat~9+,
    \texttt{jackson-mapper-asl} $\to$ \texttt{jackson-databind}~2.x &
    SBOM & A.8.8 \\
\bottomrule
\end{tabularx}
}

\section{Bewust Niet-gerealiseerde Items}
\label{sec:niet-gerealiseerd}

Conform de Sprint-4-vereiste legt deze sectie vast welke onderdelen
\emph{bewust niet} (of nog niet) zijn gemitigeerd, met per item de reden, het
restrisico en de aanbeveling voor de opvolgende sprint. De lijst bevat alléén
punten die na de verificatie van dit rapport nog open staan; de zes uitgewerkte
bevindingen B-001 t/m B-006 zijn afgerond en staan hier dus \emph{niet} tussen.

\begin{itemize}
    \item \textbf{Quick-win component-updates (\texttt{snakeyaml} $\to$ 2.x,
          \texttt{spring-web} $\to$ 5.3.34) --- A.8.8.}
          \emph{Reden:} non-breaking, maar binnen de sprinttijd niet meer
          doorgevoerd en hertest (de patches voor \texttt{netty-handler} en
          \texttt{log4j} zijn wél gerealiseerd, PR~\#105).
          \emph{Restrisico:} bekende kwetsbaarheden (o.a.\ deserialisatie van de
          RCE-klasse) blijven in het productie-artefact.
          \emph{Aanbeveling:} doorvoeren via \texttt{<dependencyManagement>} en
          een SCA-herscan als bewijsartefact toevoegen.
    \item \textbf{\texttt{c3p0} $\to$ 0.10.x en \texttt{h2} $\to$ 2.x --- A.8.8.}
          \emph{Reden:} bewust teruggedraaid --- \texttt{c3p0} 0.10.x vereist
          Java~9+ (de module bouwt op Java~8) en \texttt{h2} 2.x breekt de
          OpenMRS test-context; \texttt{h2} draait bovendien alleen in
          \texttt{test}-scope.
          \emph{Restrisico:} beperkt (geen productie-impact voor \texttt{h2}).
          \emph{Aanbeveling:} meenemen in het Java~17-/Spring~6-platformtraject.
    \item \textbf{TLS-afdwinging op deployment-niveau (TM-I4) --- A.8.24.}
          \emph{Reden:} de module dwingt zelf geen TLS af; afdwinging hoort op
          reverse proxy / servlet-container, buiten de module-grens.
          \emph{Restrisico:} HTTP Basic-Auth-credentials reizen als plaintext
          zonder gegarandeerde TLS.
          \emph{Aanbeveling:} TLS afdwingen en documenteren op de
          deployment-laag; daarna de A.8.24-rij naar \emph{geïmplementeerd}.
    \item \textbf{Restrisico's buiten de module-scope.}
          \emph{Reden:} log-injectie in OpenMRS Core
          (\texttt{HibernateContextDAO}) en versie-disclosure van Tomcat via
          \texttt{robots.txt}/\texttt{sitemap.xml} liggen in upstream-code resp.\
          de serverconfiguratie, niet in deze module.
          \emph{Restrisico:} laag tot midden (reconnaissance-enabler).
          \emph{Aanbeveling:} upstream-issue indienen resp.\ serverconfiguratie
          aanpassen bij de deployende organisatie.
    \item \textbf{Productieschakeling naar het DML-only database-account ---
          A.8.2/A.8.3.}
          \emph{Reden:} het mechanisme is aanwezig en lokaal geverifieerd, maar
          de runtime-switch kon in deze auditomgeving niet
          container-geverifieerd worden (geen Docker-daemon beschikbaar).
          \emph{Restrisico:} laag.
          \emph{Aanbeveling:} verifiëren in een omgeving met Docker-daemon.
\end{itemize}

\newpage

% ════════════════════════════════════════════════════════════════════════════
%  7. BIJLAGEN
% ════════════════════════════════════════════════════════════════════════════
% Forceer de bijlagen-lijst op een nieuwe pagina in de inhoudsopgave, zodat
% A--H niet over de paginarand van de TOC worden afgekapt.
\addtocontents{toc}{\protect\newpage}
\appendix

\chapter{CodeQL SARIF-output}
\label{app:codeql}

\textbf{SARIF} (\emph{Static Analysis Results Interchange Format}) is het
machine-leesbare JSON-formaat waarin CodeQL zijn bevindingen rapporteert. De
workflow \texttt{.github/workflows/codeql.yml} draait de query-suite
\texttt{security-extended} (zie de onderbouwing in ADR-001) bij elke push en PR
op \texttt{main}, en uploadt het SARIF-resultaat in de stap \emph{Perform CodeQL
Analysis} \emph{rechtstreeks} naar de GitHub code-scanning-backend (de
Security-tab). Het bestand wordt dus niet als losse Actions-artefact bewaard,
maar is wel volledig herleidbaar.

Vanwege de omvang (honderden \texttt{results}- en \texttt{rules}-objecten) is het
SARIF-bestand niet integraal in dit rapport opgenomen. Het is op twee manieren
herleidbaar:

\begin{itemize}
    \item \textbf{In de Security-tab:} GitHub $\to$ Security $\to$ Code scanning,
          gefilterd op \texttt{category: /language:java}
          (\url{https://github.com/GrannyGuard/webservices-rest-audit/security/code-scanning}).
    \item \textbf{Reproduceerbaar als SARIF via de API:} het ruwe SARIF van een
          specifieke analyse (hier: commit \texttt{e476489}, run
          \texttt{27280969725}, 2026-06-10) is op te halen met:
\begin{lstlisting}[basicstyle=\ttfamily\scriptsize, breaklines=true, frame=single]
# 1) zoek de analysis-id van de gewenste commit
gh api repos/GrannyGuard/webservices-rest-audit/code-scanning/analyses \
  --jq '.[] | select(.commit_sha|startswith("e476489")) | {id, category, created_at}'

# 2) download het ruwe SARIF van die analyse
gh api -H "Accept: application/sarif+json" \
  repos/GrannyGuard/webservices-rest-audit/code-scanning/analyses/<id> \
  > codeql-e476489.sarif
\end{lstlisting}
\end{itemize}

\section*{A.1\quad Uitvoering in de pipeline}

De CodeQL-scan is een verplichte status check op \texttt{main} (zie
\S\ref{app:traceability}, A.8.25). Onderstaande Actions-run toont een geslaagde
analyse van de Java-module (\emph{Initialize CodeQL} $\to$ \emph{Build with
Maven} $\to$ \emph{Perform CodeQL Analysis}); de laatste stap voert de upload
naar de Security-tab uit.

\begin{figure}[h]
    \centering
    \auditfig{../../04-code-review/Images/succesful-codecql-run.png}{0.85}
    \caption{Geslaagde CodeQL-workflow ``Analyze Java (java)'' in GitHub
             Actions --- query-suite \texttt{security-extended}, build via Maven
             (JDK 8), looptijd $\sim$3m30s.}
    \label{fig:codeql-run}
\end{figure}

\section*{A.2\quad Representatief SARIF-fragment}

Onderstaand fragment toont de structuur van het SARIF-bestand: de
\texttt{tool.driver}-metadata met één query-definitie (\texttt{rules}) en één
bijbehorende bevinding (\texttt{results}) --- hier de \texttt{java/xss}-finding
(bevinding CR-2). De overige query-definities en bevindingen volgen hetzelfde
stramien.

\begin{lstlisting}[basicstyle=\ttfamily\scriptsize, breaklines=true, frame=single]
{
  "$schema": "https://json.schemastore.org/sarif-2.1.0.json",
  "version": "2.1.0",
  "runs": [
    {
      "tool": {
        "driver": {
          "name": "CodeQL",
          "organization": "GitHub",
          "semanticVersion": "2.18.4",
          "rules": [
            {
              "id": "java/xss",
              "name": "java/xss",
              "shortDescription": { "text": "Cross-site scripting" },
              "defaultConfiguration": { "level": "error" },
              "properties": {
                "tags": [ "security", "external/cwe/cwe-079" ],
                "security-severity": "6.1"
              }
            }
            // ... overige security-extended query-definities ...
          ]
        }
      },
      "results": [
        {
          "ruleId": "java/xss",
          "level": "error",
          "message": {
            "text": "Cross-site scripting vulnerability due to a user-provided value."
          },
          "locations": [
            {
              "physicalLocation": {
                "artifactLocation": {
                  "uri": "omod/.../web/controller/SwaggerDocController.java"
                },
                "region": { "startLine": 27, "startColumn": 12, "endColumn": 47 }
              }
            }
          ],
          "partialFingerprints": { "primaryLocationLineHash": "a1b2c3d4e5f6:1" }
        }
        // ... overige results (user-controlled-bypass, log-injection, ...) ...
      ]
    }
  ]
}
\end{lstlisting}

\section*{A.3\quad Reproduceerbare alert-export}

De SARIF-bevindingen verschijnen als \emph{code scanning alerts}. Onderstaande
\texttt{gh}-commando's reproduceren de bewijslijst machine-leesbaar, gescheiden
van de SCA/Snyk-dependency-alerts (die in Hoofdstuk~\ref{ch:sbom} thuishoren):

\begin{lstlisting}[basicstyle=\ttfamily\scriptsize, breaklines=true, frame=single]
# Alle open CodeQL code-niveau security-bevindingen (CR-2..CR-6):
gh api repos/GrannyGuard/webservices-rest-audit/code-scanning/alerts --paginate \
  -q '.[] | select(.rule.id|startswith("java/")) | select(.state=="open")
      | select(.rule.security_severity_level!=null)
      | {n:.number, rule:.rule.id, sev:.rule.security_severity_level,
         loc:(.most_recent_instance.location.path+":"
              +(.most_recent_instance.location.start_line|tostring))}'

# De 4 gedismiste "false positive"-alerts (CR-1-restpunt), incl. reden:
gh api repos/GrannyGuard/webservices-rest-audit/code-scanning/alerts --paginate \
  -q '.[] | select(.rule.id=="java/log-injection") | select(.state=="dismissed")
      | {n:.number, reason:.dismissed_reason, comment:.dismissed_comment}'
\end{lstlisting}

\noindent\textbf{Momentopname} (scan \texttt{e476489}, 2026-06-10) --- de zeven
resterende open security-alerts na het sluiten van de negen
\texttt{java/log-injection}-bevindingen (zie B-006, \S\ref{sec:b006}):

\vspace{0.4em}
{\small
\begin{tabular}{@{}r l l l c@{}}
\toprule
\textbf{\#} & \textbf{Rule} & \textbf{Sev} & \textbf{Locatie} & \textbf{Bev.} \\
\midrule
1   & \texttt{java/xss}                   & high   & \texttt{SwaggerDocController.java:27}     & CR-2 \\
908 & \texttt{java/user-controlled-bypass}& high   & \texttt{AuthorizationFilter.java:113}     & CR-3 \\
4   & \texttt{java/sensitive-log}         & high   & \texttt{LayoutTemplateProvider.java:124}  & CR-4 \\
5   & \texttt{java/tainted-arithmetic}    & high   & \texttt{RequestContext.java:172}          & CR-5 \\
6   & \texttt{java/tainted-arithmetic}    & high   & \texttt{RequestContext.java:183}          & CR-5 \\
2   & \texttt{java/overly-large-range}    & medium & \texttt{ServerLogActionWrapper.java:70}   & CR-6 \\
3   & \texttt{java/overly-large-range}    & medium & \texttt{ServerLogActionWrapper.java:70}   & CR-6 \\
\bottomrule
\end{tabular}
}

\section*{A.4\quad Visuele bewijslast --- GitHub Security-tab}

De Security-tab vóór en ná de log-injection-mitigatie (B-006). De gesloten
alerts (\texttt{state: fixed}) tonen de kwantitatief afgesloten
detect-\,$\to$\,mitigeer-\,$\to$\,verifieer-cyclus: negen
\texttt{java/log-injection}-bevindingen van open naar nul.

\begin{figure}[h]
    \centering
    \auditfig{../../04-code-review/Images/bijlage-codeql-alerts-open.png}{0.95}
    \caption{Open code-scanning alerts in de Security-tab (na de switch naar
             \texttt{security-extended}: de zeven resterende
             code-niveau-bevindingen uit A.3).}
    \label{fig:codeql-open}
\end{figure}

\begin{figure}[h]
    \centering
    \auditfig{../../04-code-review/Images/bijlage-codeql-alerts-closed.png}{0.95}
    \caption{Gesloten alerts: de negen \texttt{java/log-injection}-bevindingen
             op \texttt{state: fixed} na commit \texttt{831dafb} ---
             bewijs van B-006 (\S\ref{sec:b006}).}
    \label{fig:codeql-closed}
\end{figure}

\chapter{Dependabot Alerts Log}
\label{app:dependabot}

GitHub Dependabot vormt de primaire kwetsbaarhedenstream van de module: het matcht
de afhankelijkheden uit \texttt{pom.xml} continu tegen de GitHub Advisory Database
en opent automatisch een alert per kwetsbare component (NEN-7510:2024-2 A.8.8, zie
Bijlage~\ref{app:traceability}). Onderstaand overzicht toont de openstaande alerts;
op het moment van deze opname waren dat er \textbf{292}. Dat is hoger dan de
\textbf{264} uit Hoofdstuk~\ref{ch:sbom}: de stream is \emph{live} en groeit naarmate
nieuwe advisories verschijnen --- het cijfer in de hoofdtekst is de bevroren
audit-momentopname, geen eindtotaal. Na deduplicatie herleiden deze ruwe alerts zich
tot 75 unieke CVE's (\S\ref{sec:top-dependencies}).

\begin{figure}[h]
    \centering
    \auditfig{../../04-code-review/Images/bijlage-dependabots.png}{0.95}
    \caption{Dependabot-alertoverzicht in de GitHub Security-tab (\emph{292 Open}
             op het moment van de screenshot; de bevroren audit-momentopname in de
             hoofdtekst is 264 --- zie de toelichting hierboven):
             o.a.\ de SnakeYAML-deserialisatie-RCE, Spring Framework path traversal,
             RCE in de H2 Console en Apache Struts input-validatie. Elke regel toont
             de severity, het betrokken package en de branch.}
    \label{fig:dependabot-overview}
\end{figure}

\noindent Het detail van één alert illustreert de informatie waarop het
patch-/suppress-/accept-besluit (\S\ref{sec:contextuele-risico}) berust: het
package, de \emph{affected}- versus \emph{patched}-versie, het transitieve
dependency-pad en de impactbeschrijving.

\begin{figure}[h]
    \centering
    \auditfig{../../04-code-review/Images/high-vul-dependabot-alert.png}{0.85}
    \caption{Detail van de H2 Console JNDI-RCE (\texttt{com.h2database:h2}, affected
             \texttt{< 2.0.206}). Het dependency-pad toont dat H2 \emph{transitief}
             via \texttt{org.openmrs.test:openmrs-test} binnenkomt --- en dus
             uitsluitend in \texttt{test}-scope draait. Precies daarom zakt deze
             CVSS-9.8-bevinding in de contextuele weging
             (\S\ref{sec:contextuele-risico}) naar de onderkant van de backlog: de
             component bereikt het productie-artefact aantoonbaar niet.}
    \label{fig:dependabot-detail}
\end{figure}

\chapter{CycloneDX SBOM JSON}
\label{app:sbom}

De volledige Software Bill of Materials is een machine-leesbaar
CycloneDX~1.6-bestand en is vanwege de omvang (238 componenten ---
237 afhankelijkheden plus de module zelf, $\sim$520 kB) niet integraal in dit
rapport opgenomen. Het artefact is
op twee manieren herleidbaar:

\begin{itemize}
    \item \textbf{In de repository:}
          \texttt{docs/security/03-sbom-en-cve-advies/sbom.cdx.json}.
    \item \textbf{Als CI-artefact:} de workflow \texttt{.github/workflows/sbom.yml}
          genereert het bestand bij elke push naar \texttt{main}/\texttt{develop}
          en uploadt het als artefact \texttt{sbom-cyclonedx} (90 dagen bewaard):\\
          \url{https://github.com/GrannyGuard/webservices-rest-audit/actions/workflows/sbom.yml}
\end{itemize}

Onderstaand fragment toont de structuur (metadata + één component); de overige
236 componenten volgen hetzelfde stramien.

\begin{lstlisting}[basicstyle=\ttfamily\scriptsize, breaklines=true, frame=single]
{
  "bomFormat": "CycloneDX",
  "specVersion": "1.6",
  "serialNumber": "urn:uuid:e39d869d-48d4-3a68-b00d-a60013f432b2",
  "metadata": {
    "timestamp": "2026-06-13T14:17:17Z",
    "tools": [{ "name": "cyclonedx-maven-plugin", "version": "2.9.1" }],
    "component": {
      "group": "org.openmrs.module", "name": "webservices.rest",
      "version": "3.2.0", "type": "library"
    }
  },
  "components": [
    {
      "group": "org.springframework", "name": "spring-web",
      "version": "5.3.30", "scope": "required",
      "purl": "pkg:maven/org.springframework/spring-web@5.3.30?type=jar"
    }
    // ... 236 verdere componenten ...
  ]
}
\end{lstlisting}

\chapter{PR-overzicht en Review-commentaar}
\label{app:pr}

Elke wijziging aan \texttt{main} verloopt via een pull request die zowel
\textbf{automatisch} als \textbf{handmatig} wordt beoordeeld voordat hij gemerged
mag worden (NEN-7510:2024-2 A.8.25, zie Bijlage~\ref{app:traceability}).
Onderstaande PR-tijdlijn toont beide lagen in één beeld:

\begin{itemize}
    \item \textbf{Automatisch} --- de bots \texttt{github-advanced-security}
          (CodeQL code scanning) en \texttt{sonarqubecloud} plaatsen hun
          analyseresultaat rechtstreeks als PR-commentaar; de
          SonarCloud-\emph{quality gate} (\emph{``Quality Gate passed''}) en de
          vier geslaagde status checks zijn een harde merge-voorwaarde.
    \item \textbf{Handmatig} --- een tweede teamlid wordt als reviewer
          aangewezen, beoordeelt de diff en geeft expliciet akkoord
          (\emph{``approved these changes''}) vóór de merge. De onderliggende
          commits zijn bovendien \emph{Verified} (ondertekend).
\end{itemize}

\begin{figure}[h]
    \centering
    \auditfig{../../04-code-review/Images/PR-review.png}{0.7}
    \caption{Voorbeeld van een gereviewde pull request: de geautomatiseerde
             analyse van GitHub Advanced Security (CodeQL) en SonarCloud naast de
             handmatige goedkeuring en merge door een tweede reviewer --- zowel
             machinale als menselijke controle in één keten.}
    \label{fig:pr-review}
\end{figure}

De volledige geschiedenis van alle pull requests --- open én gesloten, met de
bijbehorende review-commentaren en status checks --- is in te zien op:
\url{https://github.com/GrannyGuard/webservices-rest-audit/pulls}

\chapter{Risicomatrix en Bow-tie Diagrammen}
\label{app:risicomatrix}

Deze bijlage bevat de \emph{volledige} risico-evaluatie waarvan Hoofdstuk~\ref{ch:bevindingen}
alleen de top (score $\geq 12$) toonde: alle drieëntwintig gescoorde risico's
(\S\ref{app:matrix-full}) plus de twee bow-tie-analyses voor de best exploiteerbare
aanvalsketens SQ1 en CD1 (\S\ref{app:bowtie-sq1}--\ref{app:bowtie-cd1}). De
score-grid en de classificatiedrempels (Kans $\times$ Impact, 1--5) staan in
\S\ref{sec:risicomatrix}. Bronversie: \texttt{01-security-audit/01.md} \S4.2--4.3
en \texttt{02-secure-pipelines/02.md} \S7.2.

\section*{E.1\quad Volledige risicomatrix (23 risico's)}
\label{app:matrix-full}

De risico's komen uit vijf bronnen: \textbf{SQ} (SonarQube), \textbf{CD} (CI/CD),
\textbf{CQL} (CodeQL), \textbf{AT} (threat model) en \textbf{PT} (penetratietest).
Score $=$ Kans $\times$ Impact.

{\small
\renewcommand{\arraystretch}{1.15}
\begin{longtable}{@{}
    >{\raggedright\arraybackslash}p{1.0cm}
    >{\raggedright\arraybackslash}p{6.0cm}
    >{\centering\arraybackslash}p{0.5cm}
    >{\centering\arraybackslash}p{0.5cm}
    >{\centering\arraybackslash}p{1.1cm}
    >{\raggedright\arraybackslash}p{1.7cm}
    >{\raggedright\arraybackslash}p{2.2cm}@{}}
\toprule
\textbf{ID} & \textbf{Omschrijving} & \textbf{K} & \textbf{I} & \textbf{Score} &
\textbf{Klasse} & \textbf{KJ} \\
\midrule
\endfirsthead
\toprule
\textbf{ID} & \textbf{Omschrijving} & \textbf{K} & \textbf{I} & \textbf{Score} &
\textbf{Klasse} & \textbf{KJ} \\
\midrule
\endhead
\midrule
\multicolumn{7}{r}{\emph{vervolgt op volgende pagina}} \\
\endfoot
\bottomrule
\endlastfoot
\textbf{SQ1}  & Kwetsbare sessie-validatie (session fixation)        & 3 & 5 & \rood{15}             & \rood{Hoog}        & KJ3, KJ1, KJ2 \\
\textbf{SQ2}  & Onvolledige switch-statements (zoeklogica)           & 4 & 4 & \rood{16}             & \rood{Hoog}        & KJ2 \\
\textbf{SQ3}  & Foute testassertions (verbergen bugs)               & 5 & 3 & \rood{15}             & \rood{Hoog}        & KJ2 \\
\textbf{SQ4}  & Niet-serializeerbaar validatie-veld                 & 3 & 3 & \amber{9}             & \amber{Midden}     & KJ5 \\
\textbf{SQ5}  & Extreem hoge cognitieve complexiteit                & 5 & 2 & \amber{10}            & \amber{Midden}     & KJ5 \\
\textbf{SQ6}  & Lege methoden zonder verklaring                     & 4 & 2 & \geel{8}              & \geel{Laag-Mid.}   & --- \\
\textbf{SQ7}  & Ongeauth.\ uitlezing global properties (\texttt{/settings.form/search}) & 4 & 5 & \rood{20} & \rood{Hoog} & KJ7, KJ4 \\
\textbf{SQ8}  & Reflected XSS in Swagger debug-endpoint             & 3 & 4 & \amber{12}$^\triangle$ & \amber{Midden}    & KJ3 \\
\textbf{SQ9}  & Host-header injectie in Swagger-spec                & 3 & 3 & \amber{9}             & \amber{Midden}     & KJ3 \\
\textbf{CD1}  & Supply chain attack (dependency)                    & 3 & 5 & \rood{15}             & \rood{Hoog}        & KJ7, KJ1, KJ2 \\
\textbf{CD2}  & Lekkage van secrets                                 & 3 & 4 & \amber{12}$^\triangle$ & \amber{Midden}    & KJ7, KJ3 \\
\textbf{CD3}  & Ongeauth.\ toegang repository/pipeline              & 2 & 5 & \amber{10}            & \amber{Midden}     & KJ7 \\
\textbf{CD4}  & Foutieve beveiligingsconfiguratie                   & 3 & 3 & \amber{9}             & \amber{Midden}     & KJ7 \\
\textbf{CD5}  & Uitval CI/CD-omgeving                               & 4 & 2 & \geel{8}              & \geel{Laag-Mid.}   & KJ5 \\
\textbf{CD6}  & Malafide code contributor                           & 2 & 5 & \amber{10}$^\triangle$ & \amber{Midden}    & KJ7, KJ1, KJ2 \\
\textbf{CQL1} & Log injection (CWE-117) via REST-params             & 4 & 3 & \amber{12}            & \amber{Midden}     & KJ6 \\
\textbf{AT1}  & Onbeveiligde API-docs als reconnaissance            & 5 & 2 & \amber{10}            & \amber{Midden}     & KJ1, KJ2 \\
\textbf{AT2}  & Bulk data-exfiltratie via REST API                  & 3 & 5 & \rood{15}             & \rood{Hoog}        & KJ1, KJ2 \\
\textbf{AT3}  & Credential-aanvallen via wachtwoordherstel          & 3 & 4 & \amber{12}$^\triangle$ & \amber{Midden}    & KJ3, KJ1, KJ2 \\
\textbf{AT4}  & Gecompromitteerde derde-partij-integratie           & 3 & 4 & \amber{12}$^\triangle$ & \amber{Midden}    & KJ1, KJ2, KJ3 \\
\textbf{PT1}  & Geen rate-limiting op authenticatie-endpoints       & 4 & 5 & \rood{20}             & \rood{Hoog}       & KJ3, KJ1, KJ2 \\
\textbf{PT2}  & Gebruikersenumeratie via `/user` API & 4 & 3 & \amber{12}$^\triangle$ & \amber{Midden}   & KJ3 \\
\textbf{PT3}  & Stack trace disclosure in foutrespons               & 4 & 3 & \amber{12}$^\triangle$ & \amber{Midden}   & KJ6, KJ1, KJ2 \\
\end{longtable}
}

\noindent\small $^\triangle$ Midden-score (9--12) die de vertrouwelijkheid of
integriteit van KJ1/KJ2/KJ3 raakt met score $\geq 10$ $\to$ conform de
risicobereidheid (\S\ref{sec:risicomatrix}) geclassificeerd als \emph{onacceptabel}.
PT1 en SQ7 overschrijden die drempel al op eigen score (20). Van de drieëntwintig
risico's vallen er eenentwintig in Midden of hoger; veertien zijn onacceptabel
(zeven \rood{Hoog} plus zeven $^\triangle$-gemarkeerde Midden-risico's).\normalsize

% ── Herbruikbare opmaak voor de bow-tie-tabellen ─────────────────────────────
\newcommand{\topevent}[1]{%
    {\setlength{\fboxsep}{6pt}%
     \colorbox{avansRed}{\parbox{2.5cm}{\centering\color{white}\small
        \textbf{TOP-EVENT}\\[4pt]#1}}}}

\section*{E.2\quad Bow-tie SQ1 --- Sessie-overname (score 15)}
\label{app:bowtie-sq1}

Twee risico's scoren nominaal hoger dan SQ1 (score~15): \textbf{PT1} (geen
rate-limiting, score~20) en \textbf{SQ2} (onvolledige switch-statements, score~16).
PT1 is als bevinding B-001 uitgewerkt en maakt deel uit van de aanvalsketen die déze
bow-tie onderbouwt (geen rate-limiting vergemakkelijkt sessie-ID-brute-force).
\textbf{SQ1} is als bow-tie gekozen omdat het een aantoonbaar exploiteerbare keten kent
(sessie-fixatie via \texttt{getRequestedSessionId()} in
\texttt{AuthorizationFilter.java}, bevestigd door pentestbevinding~F-03) én de
meeste, hoogst geclassificeerde kroonjuwelen raakt (KJ3, KJ1, KJ2). De
preventieve en reactieve barrières corresponderen met de \rood{Nu}-acties uit het
geprioriteerd advies (\S\ref{sec:aanbevelingen}).

{\small
\renewcommand{\arraystretch}{1.25}
\begin{tabularx}{\textwidth}{@{}
    >{\raggedright\arraybackslash}X
    >{\centering\arraybackslash}m{3.0cm}
    >{\raggedright\arraybackslash}X@{}}
\toprule
\textbf{Oorzaken \& preventieve barrières} & \textbf{Kritieke gebeurtenis} &
\textbf{Reactieve barrières \& gevolgen} \\
\midrule
\textbf{Oorzaken}\newline
\textbullet\ Sessie-ID niet ververst na login (session fixation)\newline
\textbullet\ \texttt{getRequestedSessionId()} accepteert client-sessie-ID zonder server-side validatie\newline
\textbullet\ Geen rate-limiting $\to$ sessie-ID's brute-forcebaar\newline
\textbullet\ Geen sessie-timeout / inactiviteitsdetectie\par\smallskip
\textbf{Preventieve barrières}\newline
\textbullet\ Sessie-ID regenereren bij login (\texttt{session.invalidate()} + nieuwe sessie)\newline
\textbullet\ Server-side sessie-validatie tegen register van actieve sessies\newline
\textbullet\ Rate-limiting / brute-force-bescherming op auth-endpoints (prio 3)\newline
\textbullet\ CodeQL/SAST-scan op sessiebeheer (actief in pipeline)
&
\topevent{Aanvaller neemt een actieve sessie over $\to$ toegang tot KJ1, KJ2, KJ3}
&
\textbf{Reactieve barrières}\newline
\textbullet\ Auth-logging op INFO/WARN met remote IP + endpoint $\to$ detecteert afwijkend sessiegebruik (prio 2)\newline
\textbullet\ Forceren van logout / sessie-invalidatie bij anomalie\newline
\textbullet\ Incident-response + AVG-meldplicht (art.\ 33/34)\par\smallskip
\textbf{Gevolgen}\newline
\textbullet\ Datalek bijzondere persoonsgegevens (AVG art.\ 9) --- meldplicht + boete AP\newline
\textbullet\ Manipulatie van medische gegevens --- risico patiëntveiligheid\newline
\textbullet\ Reputatieschade / verlies van vertrouwen \\
\bottomrule
\end{tabularx}
}

\section*{E.3\quad Bow-tie CD1 --- Supply chain attack (score 15)}
\label{app:bowtie-cd1}

CD1 is het hoogst scorende CI/CD-risico en raakt dezelfde kroonjuwelen als SQ1
(KJ1, KJ2), maar via de \emph{toeleveringsketen} in plaats van de sessielaag. Beide
bow-ties samen dekken de twee hoogst scorende aanvalsroutes van het systeem.

{\small
\renewcommand{\arraystretch}{1.25}
\begin{tabularx}{\textwidth}{@{}
    >{\raggedright\arraybackslash}X
    >{\centering\arraybackslash}m{3.0cm}
    >{\raggedright\arraybackslash}X@{}}
\toprule
\textbf{Oorzaken \& preventieve barrières} & \textbf{Kritieke gebeurtenis} &
\textbf{Reactieve barrières \& gevolgen} \\
\midrule
\textbf{Oorzaken}\newline
\textbullet\ Gehackte maintainer van een legitieme package\newline
\textbullet\ Typo-squatting (vergelijkbare pakketnaam)\newline
\textbullet\ Compromise van een public registry\par\smallskip
\textbf{Preventieve barrières}\newline
\textbullet\ SCA-scans met OSV-Scanner\newline
\textbullet\ Genereren \& controleren van SBOM's (CycloneDX)\newline
\textbullet\ Pinnen van versies via SHA-hashes (GitHub Actions-stappen)
&
\topevent{Malafide dependency bereikt en draait in de CI/CD-pipeline}
&
\textbf{Reactieve barrières}\newline
\textbullet\ Pipeline stopt automatisch bij detectie kwetsbaarheid\newline
\textbullet\ Least privilege / geïsoleerde kortlevende runners\newline
\textbullet\ Monitoren \& incident response\par\smallskip
\textbf{Gevolgen}\newline
\textbullet\ Code-executie op dev/build-omgeving\newline
\textbullet\ Exfiltratie van codebase en secrets\newline
\textbullet\ Besmette releases naar eindgebruikers \\
\bottomrule
\end{tabularx}
}

\chapter{Screenshots Code Review Bevindingen}
\label{app:screenshots}

Naast CodeQL (Bijlage~\ref{app:codeql}) draait \textbf{SonarCloud} als tweede
SAST-tool. Het is ingericht als verplichte \emph{quality gate}: een pull request
waarvan de gate faalt, kan niet worden gemerged zonder een expliciet
gedocumenteerde afweging (zie \texttt{04-code-review/04.md} \S1). Onderstaand
dashboard toont de stand op de \texttt{main}-branch.

\begin{figure}[h]
    \centering
    \auditfig{../../04-code-review/Images/bijlage-sonarcloud-security.png}{0.95}
    \caption{SonarCloud-dashboard van de module: de quality gate staat op
             \emph{Failed} (3 condities niet gehaald), met een
             \emph{Security Rating E} en twee openstaande security-issues
             (50\,\% Blocker, 50\,\% High). Dit dashboard onderbouwt de
             SonarQube-bevindingen (de SQ-reeks) uit de risicomatrix
             (\S\ref{sec:risicomatrix}).}
    \label{fig:sonarcloud}
\end{figure}

\chapter{Snyk-rapport}
\label{app:snyk}

Snyk is de derde, onafhankelijke SCA-scanner naast Dependabot en OSV-Scanner. Het
draait in de CI-job \texttt{.github/workflows/sbom.yml} (\texttt{snyk test} op elke
\texttt{pom.xml}, plus \texttt{snyk monitor}, dat het hosted rapport op
\texttt{app.snyk.io} synchroon houdt) en gebruikt een \emph{eigen}
vulnerability-database. Daardoor vlagt Snyk bevindingen die de GitHub Advisory
Database --- en dus Dependabot --- nog niet kent: de 36 aanvullende, unieke CVE's
uit \S\ref{sec:top-dependencies}.

\paragraph{Momentopname.} De screenshots in deze en de voorgaande bijlagen zijn op
\emph{verschillende momenten} in het project gemaakt. Elk is daarmee een
\emph{point-in-time}-bewijs van de toenmalige stand; niet elke getoonde bevinding
is op het moment van inleveren al gemitigeerd, en de aantallen kunnen per bijlage
verschillen (vergelijk Dependabot 264 vs.\ 292, Bijlage~\ref{app:dependabot}). De
bijlagen documenteren dus het \emph{verloop} van de scanresultaten, geen enkele
eindtoestand.

\section*{G.1\quad Hosted dashboard}

Het Snyk-dashboard scant de vier Maven-manifesten van de module afzonderlijk; de
geaggregeerde telling is \textbf{21 Critical, 199 High, 120 Medium en 38 Low}. Dit
zijn \emph{ruwe} issue-tellingen per manifest --- de drie deel-poms tonen nagenoeg
identieke aantallen omdat ze grotendeels dezelfde dependency-set delen. Het in
\S\ref{sec:top-dependencies} genoemde aantal van \emph{36 unieke aanvullende} CVE's
is de gededupliceerde, gekruisvalideerde waarde, niet deze ruwe som.

\begin{figure}[h]
    \centering
    \auditfig{../../04-code-review/Images/snyk-dashboard.png}{0.98}
    \caption{Snyk Open Source-dashboard (\texttt{app.snyk.io}) voor
             \texttt{GrannyGuard/webservices-rest-audit}: vier gescande
             \texttt{pom.xml}-manifesten met geaggregeerd 21 Critical, 199 High,
             120 Medium en 38 Low issues.}
    \label{fig:snyk-dashboard}
\end{figure}

\section*{G.2\quad Voorbeeld: kritieke authenticatie-bypass in Apache Tomcat}

Onderstaand detail toont een van de kritieke bevindingen: \textbf{CVE-2026-43512}
(\texttt{CWE-287} Improper Authentication, \textbf{CVSS 9.1}) in
\texttt{org.apache.tomcat:catalina 6.0.53}. Wanneer DIGEST-authenticatie is
geconfigureerd, authenticeert de server een aanvaller die een willekeurige,
onbekende gebruikersnaam opgeeft alsnog --- een volledige authenticatie-bypass. De
component komt \emph{transitief} binnen via \texttt{org.apache.tomcat:jasper 6.0.53}
(de JSP-engine die \texttt{webservices.rest-omod-common} meebundelt). Veelzeggend:
NVD heeft deze CVE nog niet geanalyseerd, terwijl Snyk's eigen database hem al
vlagt --- precies de meerwaarde van de derde scanner.

\begin{figure}[h]
    \centering
    \auditfig{../../04-code-review/Images/snyk-specific.png}{0.98}
    \caption{Detail van CVE-2026-43512 in \texttt{org.apache.tomcat:catalina 6.0.53}
             (CVSS 9.1, Critical; Snyk-ID
             \texttt{SNYK-JAVA-ORGAPACHETOMCAT-16691234}). Snyk meldt
             \emph{``no remediation path available''}: er is geen schone versie-bump.}
    \label{fig:snyk-tomcat}
\end{figure}

\paragraph{Advies.} De oplossing is inderdaad een nieuwere dependency-versie ---
Tomcat is gepatcht in \texttt{catalina 9.0.118}, \texttt{10.1.55} en
\texttt{11.0.22}. De module draait echter nog op de Tomcat~6-lijn (6.0.53), die
sinds eind 2016 \emph{end-of-life} is, en Snyk meldt expliciet \emph{``no remediation
path available''}: een directe \texttt{<dependencyManagement>}-override is geen
schone drop-in, omdat de sprong van Tomcat~6 naar~9+ meerdere \emph{major}-versies
overspant (gewijzigde Servlet-API; vanaf Tomcat~10 bovendien de namespace-wissel
\texttt{javax}~$\rightarrow$~\texttt{jakarta}). Dit is dus geen quick-win maar een
geplande upgrade, vergelijkbaar met de \texttt{spring-web}~6-migratie uit
\S\ref{sec:contextuele-risico}. Contextueel verzachtend: de module authenticeert via
Basic-Auth in de \texttt{AuthorizationFilter}, niet via DIGEST --- de exploitvoorwaarde
(geconfigureerde DIGEST-auth) is in de standaardconfiguratie afwezig, en Snyk meldt
\emph{``no known exploit''}. Het blijft niettemin een kritieke, in te plannen upgrade.

\chapter{Traceability Matrix}
\label{app:traceability}

Deze bijlage koppelt \textbf{vijftien} NEN-7510:2024-2 controls aan een
traceerbaar bewijsartefact (repository-instelling, CI-workflow, commit of
scan-output), conform de Sprint-4-vereiste van minimaal drie controls. De rijen
voor het Normenkader (\S\ref{sec:normenkader}) zijn na de tijdens deze audit
doorgevoerde mitigaties als \emph{geïmplementeerd} getraceerd --- inclusief
\textbf{A.8.5} (rate-limiting/sessie-rotatie, code-niveau) en \textbf{A.8.9}
(settings.form-autorisatie, commit \texttt{07612ca}). Alleen \textbf{A.8.24}
blijft als \emph{gap} getraceerd: TLS-afdwinging hoort op deployment-niveau en
valt buiten de module-grens. De overige rijen (A.8.29 en de
Sprint-3-hardeningmaatregelen) zijn aanvullende controls die buiten het
oorspronkelijke Normenkader om zijn geïmplementeerd en getest. De
onderhouden bronversie van deze matrix staat in de repository:
\texttt{docs/security/traceability-matrix.md}. Sectieverwijzingen
(\S) verwijzen naar de onderliggende documentatie in \texttt{docs/security/}.

{\small
\begin{longtable}{@{}
    >{\raggedright\arraybackslash}p{2.0cm}
    >{\raggedright\arraybackslash}p{4.3cm}
    >{\raggedright\arraybackslash}p{6.2cm}
    >{\raggedright\arraybackslash}p{2.0cm}@{}}
\toprule
\textbf{Control} & \textbf{Maatregel en implementatie} &
\textbf{Bewijsartefact} & \textbf{Status} \\
\midrule
\endfirsthead
\toprule
\textbf{Control} & \textbf{Maatregel en implementatie} &
\textbf{Bewijsartefact} & \textbf{Status} \\
\midrule
\endhead
\midrule
\multicolumn{4}{r}{\emph{vervolgt op volgende pagina}} \\
\endfoot
\bottomrule
\endlastfoot

\textbf{A.8.3} Toegangs\-beveiliging &
Branch protection op \texttt{main} (1 verplichte review, geen directe
pushes); GitHub Environment \texttt{prod} dwingt goedkeuring van één van twee
aangewezen reviewers af vóór de release-pipeline draait (\texttt{environment:
prod}), naast een protected-branch-policy. &
\texttt{.github/workflows/release-sign.yml} (\texttt{environment: prod});
\texttt{gh api repos/GrannyGuard/webservices-rest-audit/environments/prod}
(geverifieerd 2026-06-10); \texttt{02-secure-pipelines/02.md} \S2 en \S4.1. &
Geïmplementeerd \\
\addlinespace

\textbf{A.8.5} Authenticatie (CI/CD) &
Org-brede 2FA-verplichting voor alle leden; secret scanning met push
protection op de repository. &
\texttt{gh api orgs/GrannyGuard} toont
\texttt{two\_factor\_requirement\_enabled: true} (geverifieerd 2026-06-10);
0 open secret-scanning-alerts; \texttt{02.md} \S4.2--4.3. &
Geïmplementeerd \\
\addlinespace

\textbf{A.8.5} Authenticatie (code) &
Brute-force-bescherming en sessiehardening in \texttt{AuthorizationFilter}:
in-memory IP-rate-limiter (max 10 mislukte pogingen/60s $\to$ HTTP 429),
sessie-rotatie na \texttt{Context.authenticate()} tegen session fixation, en
privilege-check + uniforme \texttt{AUTH\_FAILURE}-logging tegen
gebruikersenumeratie. &
Commits \texttt{59bf421} (rate-limiter), \texttt{9c2fb64} (sessie-rotatie),
\texttt{ed6f3da} (enumeratie); tests \texttt{AuthorizationFilterRateLimitTest}
(6) + \texttt{AuthorizationFilterSessionFixationTest} (5); pentest-hervalidatie
\texttt{05-penetration-tests} \S18. &
Geïmplementeerd + hertest \\
\addlinespace

\textbf{A.8.8} Beheer van technische kwetsbaarheden &
SBOM (CycloneDX) bij elke push; SCA via Dependabot, OSV-Scanner en Snyk;
contextuele risicoscore met patch-/suppress-/accept-besluit per
kwetsbaarheid; updates handmatig (geen automatische Dependabot-PR's, zie
commit \texttt{e724ba1}). &
\texttt{.github/workflows/sbom.yml}; \texttt{sbom.cdx.json}
(238 componenten, zie Bijlage~\ref{app:sbom}); 264 open Dependabot-alerts
herleid tot 75 unieke CVE's; geprioriteerde backlog in \texttt{03.md}
\S4--5. &
Geïmplementeerd \\
\addlinespace

\textbf{A.8.15} Logging &
CodeQL \texttt{security-extended} als security-SAST; CWE-117
log-injection-bevindingen gemitigeerd via centrale
\texttt{RestUtil.sanitizeForLog()}. &
\texttt{.github/workflows/codeql.yml} (query-suite-keuze in ADR-001); fix
in commit \texttt{831dafb}; CodeQL-run \texttt{27280969725} op commit
\texttt{e476489} (2026-06-10): alle 9 \texttt{java/log-injection}-alerts
(\#7--15) op \emph{fixed}; zie Bijlage~\ref{app:codeql}. &
Geïmplementeerd en geverifieerd \\
\addlinespace

\textbf{A.8.15} Logging (auth-events) &
Zes logging-gaps gedicht conform NEN-7510~8.15: auth-succes en -falen
verhoogd van DEBUG naar INFO/WARN met gebruikersnaam, remote IP en URI;
IP-blokkering (\texttt{ACCESS\_BLOCKED}), sessie-timeout
(\texttt{SESSION\_TIMEOUT}), logout (\texttt{AUTH\_LOGOUT}) en toegang tot
het diagnose-endpoint (\texttt{DIAG\_ACCESS}) worden nu gelogd.
Wachtwoorden, sessie-ID's en rollen worden expliciet \emph{niet}
gelogd. &
Commit \texttt{8521cdc}; 16 geautomatiseerde tests in
\texttt{AuthorizationFilterLoggingTest} en
\texttt{SessionController1\_9LoggingTest} slagen (\texttt{mvn test -pl
omod-common,omod}); zie \texttt{01-security-audit/01.md} \S3.3. &
Geïmplementeerd en geverifieerd \\
\addlinespace

\textbf{A.8.25} Veilige ontwikkelcyclus &
Wijzigingen aan \texttt{main} uitsluitend via PR met verplichte review;
CodeQL als required status check blokkeert merges met security-findings. &
GitHub Ruleset \emph{main protection}
(\texttt{required\_approving\_review\_count: 1}; required status check
\emph{Analyze Java (java)}); \texttt{02.md} \S4.1 en \S8. &
Geïmplementeerd \\
\addlinespace

\textbf{A.8.9} Configuratiebeheer &
\texttt{settings.form/search} retourneerde alle global properties (incl.\
\texttt{REST\_ALLOWED\_IPS}) zonder authenticatie --- SQ7, risicoscore 20,
hoogste van de analyse. Nu afgeschermd met \texttt{Context.isAuthenticated()}
+ \emph{Manage RESTWS}-privilege vóór uitlezing. &
Commit \texttt{07612ca}; test \texttt{SettingsFormControllerTest} (unauth
geweigerd / geautoriseerde admin toegestaan);
\texttt{01-security-audit/01.md} \S3.5. &
Geïmplementeerd + getest \\
\addlinespace

\textbf{A.8.24} Cryptografie --- \emph{gap} &
Module dwingt zelf geen TLS af; HTTP Basic Auth-credentials reizen als
plaintext zonder gegarandeerde TLS op deployment-niveau (TM-I4). CI/CD-secrets
zijn wel beschermd (secret scanning, zie A.8.5). &
\texttt{01-security-audit/01.md} \S3.4 (gap-analyse);
\texttt{AuthorizationFilter.java:86--114}. &
\amber{Gedeeltelijk} \\
\addlinespace

\textbf{A.8.29} Beveiligings\-tests &
Pentest op VirtualBox-testomgeving (Kali Linux 2024, OpenMRS~2.8.7, OWASP WSTG v4.2):
ZAP-scan + handmatige tests; 9 bevindingen (F-01 t/m F-09), 1~kritisch (geen
rate limiting, credentials achterhaald in 6~pogingen), 4~hoog, 2~medium, 2~laag;
gekoppeld aan SQ*/TM-*-risico's en B-001 t/m B-006. \textbf{Hervalidatie
(post-mitigatie) afgerond} (Kali 2026-06-16/17): F-01 t/m F-09 hertest,
mitigaties bevestigd. &
Bijlage~\ref{app:pentest}; \texttt{05-penetration-tests/\allowbreak PentestReport-OpenMRS.md}
(F1--F9 + \S18 Hervalidatie); \texttt{zap-report.html}. &
Geïmplementeerd en geverifieerd \\
\addlinespace

\textbf{A.8.2 / A.8.3} Least-privilege database &
Twee-account-DB-model: privileged migratie-account (Liquibase/DDL)
gescheiden van een DML-only runtime-account
(\texttt{SELECT/INSERT/UPDATE/DELETE}, kan nooit \texttt{DROP}). &
\texttt{docker/\allowbreak db-init/\allowbreak 10-runtime-least-privilege.sh},
\texttt{docker-compose.yml}, \texttt{docker/\allowbreak prod/\allowbreak docker-compose.yml};
hardening checklist \S2. &
Mechanisme (PoC)$^\ast$ \\
\addlinespace

\textbf{A.8.20 / A.8.26} Host-validatie (Swagger/OpenAPI) &
OpenAPI-spec reflecteert \texttt{Host}/\texttt{X-Forwarded-Proto} niet langer
ongevalideerd: Host-allow-list + scheme-sanitisatie (SQ9); Swagger-endpoints
uitschakelbaar via \texttt{enableSwaggerDocs}. &
\texttt{RestUtil.\allowbreak resolveSwaggerHost()}/\allowbreak
\texttt{sanitizeScheme()}/\allowbreak
\texttt{isSwaggerDocsEnabled()}; unit tests
\texttt{SwaggerHardeningRestUtilTest} (8/8 groen); hardening checklist
\S1, \S10. &
Geïmplementeerd en getest \\
\addlinespace

\textbf{A.5.23 / A.8.32} Supply chain security &
Reproduceerbare, verifieerbare releases: keyless cosign-signing + SLSA
build-provenance + CycloneDX-SBOM naast de \texttt{.omod}; alle CI-actions
SHA-pinned (was 1/9). &
\texttt{.github/\allowbreak workflows/\allowbreak release-sign.yml}; SHA-pins in alle workflows;
hardening checklist \S8, \S9. &
Geïmplementeerd (workflow) \\
\addlinespace

\textbf{A.8.28} Veilig coderen &
Server-side validatie i.p.v.\ client-side: Swagger Host-allow-list +
scheme-sanitisatie (SQ9), reflected-XSS output-encoding
(\texttt{HtmlUtils.htmlEscape()} in \texttt{SwaggerDocController.debug()}, SQ8)
en centrale CWE-117 log-injectie-neutralisatie
(\texttt{RestUtil.sanitizeForLog()}, CQL1). &
\texttt{RestUtil.java}, \texttt{SwaggerDocController.java}; tests
\texttt{SwaggerHardeningRestUtilTest} (8/8) + \texttt{RestUtilTest}; CWE-79-fix
commit \texttt{e3f3c7d}, CWE-117-fix commit \texttt{831dafb}. &
Geïmplementeerd en getest \\
\addlinespace

\textbf{A.8.31} Scheiding OTAP-omgevingen &
Aparte compose-overrides per omgeving + GitHub Environments
(dev/test/prod) met gescheiden secrets; prod publiceert de DB niet en faalt
fast (\texttt{:?}) op ontbrekende secrets. &
\texttt{docker/\allowbreak\{dev,test,prod\}/\allowbreak docker-compose.yml},
\texttt{docker/\allowbreak README.md}; \texttt{02.md} \S2. &
Geïmplementeerd \\

\end{longtable}
}

$^\ast$ Het mechanisme is aanwezig en lokaal geverifieerd; de
prod-runtime-switch naar het DML-only account is een deploy-stap die in deze
omgeving niet container-geverifieerd kon worden (Docker-daemon niet
beschikbaar).

\section*{Restpunten}

\begin{itemize}
    \item Voor A.8.3 en A.8.5: screenshots toevoegen als visueel bewijs
          naast de \texttt{gh api}-output
          (GitHub Environment-instellingen en org-2FA-policy).
    \item A.8.8 aanvullen met PR en hernieuwde scanresultaten zodra de twee
          quick-win-patches (\texttt{snakeyaml} naar 2.x,
          \texttt{spring-web} naar 5.3.34) zijn doorgevoerd (de
          \texttt{netty-handler}- en \texttt{log4j}-pins zijn reeds gerealiseerd,
          PR~\#105).
    \item A.8.24 (\emph{gap}-rij) bijwerken naar ``Geïmplementeerd'' zodra de
          TLS-afdwinging op deployment-niveau is gerealiseerd. \emph{(A.8.9 is
          inmiddels gemitigeerd, commit \texttt{07612ca}.)}
    \item Voor A.8.2/A.8.3 (least-privilege database): de prod-runtime-switch
          ($^\ast$) container-verifiëren zodra een omgeving met
          Docker-daemon beschikbaar is.
\end{itemize}

\chapter{Penetratietest Samenvatting}
\label{app:pentest}

Deze bijlage bevat de samenvatting van de penetratietest die is uitgevoerd op de
OpenMRS REST Webservices module v3.2.0. De volledige rapportage bevindt zich in de
repository onder \texttt{docs/security/05-penetration-tests/PentestReport-OpenMRS.md}.

\section*{I.1\quad Testomgeving en scope}

\begin{tabular}{@{}ll@{}}
\toprule
\textbf{Component} & \textbf{Waarde} \\
\midrule
Doelstelling & OpenMRS REST Webservices module v3.2.0 \\
Server & Apache Tomcat 9.0.118 \\
Database & MariaDB 10.11.18 \\
OpenMRS versie & 2.8.7 \\
Testplatform & Kali Linux 2024 (VirtualBox, geïsoleerd) \\
Testperiode & 11--12 juni 2026 \\
Methodologie & OWASP Web Security Testing Guide v4.2 \\
Autorisatie & Schoolopdracht LU2 Groep 17, geen externe systemen bereikt \\
\bottomrule
\end{tabular}

\medskip

\textbf{In scope:} \texttt{/openmrs/ws/rest/v1/*} --- authenticatie, autorisatie,
injectie, sessiebeheer, serverconfiguratie.
\textbf{Buiten scope:} OpenMRS webinterface, directe database-toegang,
netwerkaanvallen, social engineering.

\section*{I.2\quad Bevindingen overzicht}

{\small
\renewcommand{\arraystretch}{1.15}
\begin{longtable}{@{}
    >{\raggedright\arraybackslash}p{0.9cm}
    >{\raggedright\arraybackslash}p{4.7cm}
    >{\raggedright\arraybackslash}p{1.8cm}
    >{\raggedright\arraybackslash}p{1.6cm}
    >{\raggedright\arraybackslash}p{1.8cm}
    >{\raggedright\arraybackslash}p{2.3cm}@{}}
\toprule
\textbf{ID} & \textbf{Bevinding} & \textbf{CWE} &
\textbf{Ernst} & \textbf{NEN-7510} & \textbf{Status} \\
\midrule
\endfirsthead
\toprule
\textbf{ID} & \textbf{Bevinding} & \textbf{CWE} &
\textbf{Ernst} & \textbf{NEN-7510} & \textbf{Status} \\
\midrule
\endhead
\midrule\multicolumn{6}{r}{\emph{vervolgt}}\\\endfoot
\bottomrule\endlastfoot

F-01 & Geen rate limiting op authenticatie &
    CWE-307 & \rood{Kritiek} & A.8.5 & \groen{Gemitigeerd} \\
F-02 & Gebruikersenumeratie via `/user` API &
    CWE-204 & \amber{Hoog} & A.8.5 & \groen{Gemitigeerd} \\
F-03 & Session fixation --- JSESSIONID niet geroteerd na login &
    CWE-384 & \amber{Hoog} & A.8.3 & \groen{Gemitigeerd} \\
F-04 & Stack trace disclosure in foutrespons &
    CWE-209 & \amber{Hoog} & A.8.15 & \groen{Gemitigeerd} \\
F-05 & Versie-informatie via \texttt{/systeminformation} &
    CWE-200 & \amber{Hoog} & A.8.15 & \groen{Gemitigeerd} \\
F-06 & Integer overflow in pagineringsparameter &
    CWE-190 & \geel{Middel} & A.8.28 & \groen{Gemitigeerd} \\
F-07 & Log injection via Basic Auth-gebruikersnaam &
    CWE-117 & \geel{Middel} & A.8.15 & \groen{Gemitigeerd} \\
F-08 & Ontbrekende security headers &
    CWE-693 & Laag & A.8.23 & \groen{Gemitigeerd} \\
F-09 & Cookie zonder SameSite-attribuut &
    CWE-1275 & Laag & A.8.23 & \groen{Gemitigeerd} \\

\end{longtable}
}

\section*{I.3\quad Meest kritieke bevinding --- F-01}

\textbf{Geen rate limiting op authenticatie (CWE-307, KRITIEK).}
Het authenticatie-endpoint \texttt{GET /openmrs/ws/rest/v1/session} accepteert
onbeperkt Basic Auth-pogingen zonder vertraging, blokkade of HTTP~429 op
netwerkniveau. OpenMRS heeft wel een per-account lockout voor \emph{bestaande}
gebruikers (na circa 7 pogingen), maar niet voor niet-bestaande gebruikersnamen.
Tijdens de test zijn de geldige admin-credentials
(\texttt{admin:Admin123!}) verkregen in poging~6 van~84, binnen \textbf{2 seconden},
door een Python-script en Hydra --- ruim vóór de account-lockout sloeg.

Dit levert volledige administratieve toegang op tot alle patiëntgegevens. In een
zorgomgeving is dit een AVG-meldingswaardige inbreuk (NEN-7510 A.8.5). Aanbeveling:
rate limiting op HTTP-niveau (max.\ 10 pogingen/min per IP, HTTP 429), exponentiële
vertraging bij herhaalde fouten, en alerting op brute force-patronen.

\section*{I.4\quad Mitigatiestatus}

Alle 9 bevindingen (F-01 t/m F-09) zijn gemitigeerd en op de testomgeving
(Kali, 2026-06-16/17) hertest. Per finding de fix-commit:

{\small
\begin{tabular}{@{}l l l@{}}
\toprule
\textbf{Finding} & \textbf{Commit} & \textbf{Hertest (Kali)} \\
\midrule
F-01 (rate limiting, B-001)      & \texttt{59bf421} & HTTP 429 op 11e poging \\
F-02 (gebruikersenumeratie, B-003) & \texttt{ed6f3da} & privilege-blokkade + uniforme logs \\
F-03 (session fixation, B-002)   & \texttt{9c2fb64} & JSESSIONID geroteerd na login \\
F-04 (stack trace, B-005)        & \texttt{9c30cdd} & leeg \texttt{detail}-veld + HTTP 403 \\
F-05 (versie-informatie)         & \texttt{ef30066} & DB-credentials afwezig, 401 ongeauth. \\
F-06 (integer overflow)          & \texttt{d185b31} & HTTP 500 i.p.v.\ 200 bij overflow \\
F-07 (log injection, B-006)      & \texttt{831dafb} & newlines gestript (CodeQL 9 \textrightarrow{} 0) \\
F-08 (security headers)          & \texttt{077d2b1} & alle 6 headers aanwezig \\
F-09 (cookie SameSite)           & \texttt{2d496a6} & \texttt{SameSite=Strict} op JSESSIONID \\
\bottomrule
\end{tabular}
}

\vspace{0.4em}
\noindent Twee restrisico's blijven \emph{buiten de module-scope} bestaan en zijn
als zodanig gedocumenteerd (zie \S\ref{sec:niet-gerealiseerd}): log-injectie in
OpenMRS Core (\texttt{HibernateContextDAO}) en versie-disclosure van Tomcat via
\texttt{robots.txt}/\texttt{sitemap.xml}. Voor F-07 is de mitigatie binnen de
module (de \texttt{AuthorizationFilter}) dus volledig; het OpenMRS-Core-deel
vereist een upstream-fix.

\end{document}
